% Stage G: the delivery of the points to the callback of the caller, with the stop that it may
% request; the options of the output RATPOINTS_NO_CHECK and RATPOINTS_NO_Y; the flags written back;
% the order of the points stated in terms of the polynomial; and the options of the output, the
% callback and the report of the command-line program. Proposed names are set with \lname and
% become \lean once the declarations exist.
\chapter{The delivery of the points and the output}\label{chap:output}

Stage~G models how ratpoints hands its points to its caller. The library function \verb|find_points|
returns no list: for each point that it finds it calls a function \verb|process| of the caller with
the coordinates, the caller's data \verb|info| and a pointer to a flag \verb|quit|, adds the value
that the call returns to a count \verb|total|, which it returns, and stops the search as soon as a
call sets the flag. Until now the model returned the list of the points, and the program printed it.
G threads the callback and a \emph{delivery}, the data, the count and the flag, through the
components that the C code passes \verb|process| and \verb|info| through, and tests the flag where
the C code tests \verb|quit|: in the exact check and at infinity between the calls with $+y$ and
$-y$, in the walk through the set bits of a word after each check, after each piece of bit arrays,
after each denominator, and after the points at infinity, where the search returns before it sieves.
The search with a callback is \verb|findPointsWith|; \verb|findPoints| keeps its type and becomes
the instance of it whose callback appends each point to an array and never stops the search, so
that, with the two options below clear, it returns the list of before. With the delivery come the
two options of the library that change what it delivers: \verb|RATPOINTS_NO_Y|, under which only the
point with $y \ge 0$ of a pair goes to the callback, and \verb|RATPOINTS_NO_CHECK|, under which
every numerator that passes the sieve goes to it with $y = 0$, without the exact check. They become
the fields $\mathrm{noCheck}$ and $\mathrm{noY}$ of the options of the search, bits of the word of
flags that also holds the switch of the reversal; and the bits of the flags that the program reads
after the search, whether the polynomial was reversed and whether the denominators were the squares
or divisors of the leading coefficient times squares, are written back in the result.

G also states the order of the points in terms of the polynomial, as Chapter~\ref{chap:spec}
promised: the specification gains a file that defines the list of the points in the order of
ratpoints, the points at infinity of the polynomial searched first, then its denominators in the
order of the loop of the search, ascending but for the divisors of the leading coefficient in the
order of binary counting over its primes, the numerators ascending and $+y$ before $-y$, mapped back
after a reversal; and a theorem says that \verb|findPoints| returns that list. And G gives the
command-line program the options of the output of \verb|main.c| (\verb|-x -X -y -Y -z -Z -1 -i -I|
and the formats and strings of \verb|-f -fs -fm -fe|), its callback \verb|process|, which prints the
points as they are delivered, and the report that \verb|main.c| prints without \verb|-q|, but for
the messages on the search region and the search intervals (stage~H) and on the moduli used
(stage~I); a theorem says what the program prints, in particular that \verb|-1| prints the first
point of the list that \verb|-i| keeps.

With the appending callback the delivery of a list is that list, so \verb|findPoints| returns the
list of before G whenever neither option of the output is set, and \verb|Params| and
\verb|Params.WF| do not change. Of the statements that comparator certifies, membership gains the
hypothesis that the exact check is on and, under \verb|RATPOINTS_NO_Y|, the conclusion $y \ge 0$,
since under \verb|RATPOINTS_NO_CHECK| the list holds survivors that are no points and under
\verb|RATPOINTS_NO_Y| only one point of each pair; no duplicates keeps its text, which now
quantifies over the options of the output too; the two theorems of the reference, both refusal
theorems and the well-formedness of the parameters of the program keep their text. Five statements
are new: the order theorem, over the file of the order, which the challenge restates verbatim; the
protocol of the delivery, by which the search with any callback gives the data and the count of the
delivery of the list of \verb|findPoints|, up to the first stop; and three statements on the output
under \verb|RATPOINTS_NO_CHECK|. So comparator certifies twelve statements, each brought in with its
theorem.

As in the other stages, since the aim of the project is to verify the C code as far as that can be
done, the model follows the C code where the C code has its data and its tests, at its interface
too: the callback takes the C's arguments $(x, z, y)$ and the caller's data and gives back the
weight, the stop and the new data; the flag is tested where the C code tests \verb|quit|; the points
at infinity are delivered as the C code delivers them; and the bits of the flags that the program
reads are written back as \verb|find_points| writes them. So a stop stops the search:
\verb|ratpoints-lean '1 0 1' 1000000000000 -q -1| returns at once after the point at infinity, as
ratpoints does. One thing still differs until the blocks of stage~I: the model lists the
denominators of a shape before it sifts the first one, all of them, also those after a stop. The
speed of the model is not an aim.

\section{What the C code does}\label{sec:outputC}

\paragraph{The callback.} The prototypes of \verb|find_points| and \verb|find_points_work| give the
callback as \verb|int proc(long x, long z, const mpz_t y, void *info, int *quit)|
(\csource{ratpoints.h}{436}{437}, \csource{ratpoints.h}{441}{442}). The library passes \verb|info|
through unchanged, adds the value that each call returns to its count \verb|total|, which the manual
calls a weight for counting the points, ``usually 1'', and returns the count, or one of the negative
codes of an error (\csource{ratpoints.h}{427}{430}, with the comment of \verb|RATPOINTS_NO_MEMORY|
up to line 433). Every refusal of the input (\verb|RATPOINTS_NON_SQUAREFREE|,
\verb|RATPOINTS_BAD_ARGS|, \verb|RATPOINTS_WORK_LENGTH_TOO_SMALL|) comes before the first call.
\verb|RATPOINTS_NO_MEMORY| can come after calls: once the points at infinity are delivered, the
search can still fail for lack of memory for its arrays (lines 5127-5128 and 5137-5139 of
\verb|find_points.c|), or for the points of a block on a sieving thread (lines 5146-5148); the
points delivered before it are points of the curve, but the search is incomplete, as the manual
says. The model cannot meet this code, since its search has no allocation that can fail, nor
\verb|RATPOINTS_WORK_LENGTH_TOO_SMALL|. A call stops the search by setting \verb|*quit|, which the
library clears before the search and tests after each call, so that it is clear at every call. With
threads (stage~I) the calling thread makes the same calls in the same order. The word
\verb|args->flags| holds the options of the caller (\csource{ratpoints.h}{403}{411}): without
\verb|RATPOINTS_NO_CHECK|, \verb|RATPOINTS_NO_Y| is the only one that changes which points are
delivered. \verb|RATPOINTS_NO_REVERSE| changes only the order of the calls: the search runs in the
reversed coordinates and maps each point back. \verb|RATPOINTS_NO_JACOBI| changes nothing that is
delivered. Under \verb|RATPOINTS_NO_CHECK|, which delivers every survivor of the sieve with $y = 0$,
both of them also change which pairs are delivered (see the end of the next paragraph).
\verb|RATPOINTS_VERBOSE| changes only what the library prints; \verb|find_points_work_1| keeps the
bits of the input (\csource{find_points.c}{4606}{4606}) and sets the internal ones
(\csource{ratpoints.h}{414}{421}).

\paragraph{The calls.} There are points at infinity when the degree $d$ of the polynomial searched
is odd or $c_d$ is a square (\verb|point_at_infty|, \csource{find_points.c}{4894}{4899}); then
\verb|find_points_work_1| calls the callback with $(a, b) = (1, 0)$, or $(0, 1)$ after a reversal
(\csource{find_points.c}{5068}{5099}): under \verb|RATPOINTS_NO_CHECK| once, with $y = 0$; otherwise
for odd $d$ once, with $y = 0$, and for even $d$ with $y = \sqrt{c_d}$ and then, unless that call
set \verb|quit| or \verb|RATPOINTS_NO_Y| is set, with $-\sqrt{c_d}$. If a call set \verb|quit|, it
returns the count (\csource{find_points.c}{5100}{5103}) and does not sieve: \verb|sieving_info| has
run before the points at infinity, and the block of the sieving
(\csource{find_points.c}{5115}{5167}) comes after. In the sieve, each numerator $a$ of the
denominator $b$ that passes the test of a candidate goes through \verb|RP_CHECK_POINT|
(\csource{sift.c}{202}{205}) to \verb|_ratpoints_check_point| (\csource{find_points.c}{5250}{5338}):
without \verb|RATPOINTS_NO_CHECK| it computes $v = F(a, b)$, times $b$ for odd degree, and $y =
\lfloor \sqrt{v} \rfloor$, and when $v \ge 0$ and $v = y^2$ it calls the callback with the point
mapped back after a reversal, $(b, y, a)$ for $a \ge 0$ and $(-b, y, -a)$ for $a < 0$, and then,
unless that call set \verb|quit|, $y = 0$ or \verb|RATPOINTS_NO_Y| is set, with $-y$; under
\verb|RATPOINTS_NO_CHECK| it calls the callback once, with the pair mapped in the same way and $y =
0$, whether or not it gives a point. \verb|RATPOINTS_NO_CHECK| is read nowhere else but in a line of
the output under \verb|RATPOINTS_VERBOSE|, \verb|RATPOINTS_NO_Y| nowhere else at all: the reversal,
the Sturm step, the exits, the tests on the denominators, the sieve and the third stage run as
without them. Conversely, the moduli, the third stage, the Jacobi test, the reversal and the search
region of the Sturm step change what \verb|RATPOINTS_NO_CHECK| delivers: the numerators that pass
the sieve.

\paragraph{The stop.} The flag is tested only right after a call of the callback, and each function
returns its count at once when it is set: \verb|_ratpoints_sift0| after each \verb|RP_CHECK_POINT|
(\csource{sift.c}{1250}{1251} and \csource{sift.c}{1287}{1288}, the first word of a bit array and
the others, and \csource{sift.c}{1168}{1169} in the arm of \verb|USE_LONG_IN_PHASE_2|; the arms
under \verb|DEBUG|, lines 1240-1241 and 1279-1280, are the same walk without their output,
Definition~\ref{def:walkWord}); \verb|sift| after each call of \verb|_ratpoints_sift0|
(\csource{find_points.c}{3611}{3611}; line 3610 is the stop of a thread of the pool, stage~I);
\verb|run_block| after each \verb|sift|, in each shape (\csource{find_points.c}{3928}{3928},
\csource{find_points.c}{3957}{3957}, \csource{find_points.c}{4044}{4044},
\csource{find_points.c}{4068}{4068}); and the loop over the blocks after each block
(\csource{find_points.c}{5158}{5158}). The points at infinity use a local flag
(\csource{find_points.c}{4571}{4571}), the sieve the flag of the sieving thread, which
\verb|worker_init| clears (\csource{find_points.c}{3129}{3129}); the first is clear when the sieve
starts, since a stop at infinity returns before it. So the calls made are those of the search that
no call stops, up to and including the first call that sets \verb|quit|, and the count is the sum of
the values that these calls return.

\paragraph{The flags written back.} \verb|find_points_work_1| sets \verb|RATPOINTS_REVERSED| in the
reversal check and \verb|RATPOINTS_USE_SQUARES| and \verb|RATPOINTS_USE_SQUARES1|
(\csource{find_points.c}{4910}{4943}) before both exits and before the points at infinity, so that
they are written back on every path that does not refuse the input; the return of the Sturm step at
line 4889 (stage~H) comes for even degree only, where both are clear. The program reads them after
the search, with \verb|RATPOINTS_NO_JACOBI| and \verb|RATPOINTS_NO_CHECK|, which it set itself
(message 7). The numbers of the moduli used, \verb|sp1_used|, \verb|sp2_used| and \verb|sp3_used|
(\csource{ratpoints.h}{378}{378}), which message 4 reads, belong to the correction of the moduli
(stage~I). The search intervals, \verb|num_inter| and \verb|domain|, which the search leaves as it
made them (lines 4831-4841 and the Sturm step) and message 8 reads, belong to the search region
(stage~H).

\paragraph{The program.} \verb|main.c| keeps the settings of the output in global variables
(\csource{main.c}{57}{64}), with the defaults of a command line without options
(\csource{main.c}{383}{390}), and tells an option by the character after its \verb|-|, and \verb|-f|
by the one after that: \verb|-x| and \verb|-X| set and clear \verb|RATPOINTS_NO_CHECK|, \verb|-y|
and \verb|-Y| \verb|RATPOINTS_NO_Y|, \verb|-z| and \verb|-Z| whether no point is printed, \verb|-1|
that one point is enough, and \verb|-i| and \verb|-I| whether the points at infinity are passed over
(\csource{main.c}{602}{637}); \verb|-f|, \verb|-fs|, \verb|-fm| and \verb|-fe| take the next word,
whatever it is, as the format of a point and the strings before, between and after the points
(\csource{main.c}{533}{577}), a missing word and \verb|-f| followed by another character being error
6; the last setting wins. Without \verb|-f| the format is \verb|(%x : %z)\n| under
\verb|RATPOINTS_NO_CHECK| or \verb|RATPOINTS_NO_Y| and \verb|(%x : %y : %z)\n| otherwise
(\csource{main.c}{681}{685}), and \verb|-q| clears \verb|RATPOINTS_VERBOSE|
(\csource{main.c}{688}{688}). The data that \verb|main.c| gives \verb|process|
(\csource{main.c}{110}{113}, filled at \csource{main.c}{235}{242}) holds whether a point was
printed, the settings it reads, and \verb|no_y|, set under either option of the output.
\verb|process| (\csource{main.c}{139}{189}) passes a point at infinity over under \verb|-i| with the
weight $0$, leaving \verb|quit| clear; it passes a point outside the search region over in the same
way (lines 150-152, stage~H); otherwise it sets \verb|*quit| to the setting of \verb|-1|, returns
$1$ under \verb|-z|, and else prints the string before the points if it printed no point yet and the
string between them otherwise, then the format character by character, \verb|%x|, \verb|%y| and
\verb|%z| as the coordinates (\verb|%y| as it stands under \verb|no_y|), a backslash with \verb|t|,
\verb|n|, a backslash or \verb|%| as a tab, a line feed, a backslash or a percent sign, a trailing
\verb|%| or backslash as itself, every other pair of characters as it is, and returns $1$. After the
search \verb|main.c| prints the string before the points if \verb|process| printed none, then the
string after them (\csource{main.c}{265}{265}, \csource{main.c}{267}{267}); unless \verb|-q| it
prints the report: before the search the banner, the curve (\verb|print_poly|,
\csource{main.c}{913}{949}) and, unless \verb|-v|, the height bound and a blank line
(\csource{main.c}{215}{220}), with the search region of the command line (line 224, stage~H); after
it, unless \verb|-v|, two line feeds and the messages 4 on the moduli used, 7 on what the search did
and 8 on the search intervals, then message 2 with the count, ``rational points found'' or, under
either option of the output, ``rational point pairs found'' (\csource{main.c}{268}{276}, the
messages \csource{main.c}{972}{1033}).

\paragraph{The order.} The calls of a search that no call stops come in this order, for the
polynomial that the reversal check hands on: the points at infinity, then, for the denominators in
the order of the loop of the shape, the numerators in increasing order, $+y$ before $-y$, each point
mapped back after a reversal. The denominators ascend, except for odd $d$ with $|c_d| \ne 1$ when
\verb|setup_us1| succeeds: when its factorization of $|c_d|$ by trial division, with $2$ and the odd
primes of the table up to $1021$ and a cofactor below $p^2$ taken as a prime, is complete
(\csource{find_points.c}{777}{860}, the test at \csource{find_points.c}{863}{863}), and the
squarefree divisors up to $b_{\mathrm{high}}$, made by appending $x p$ for each prime $p$ in turn to
every divisor $x \le \lfloor b_{\mathrm{high}}/p \rfloor$ made so far, are fewer than
\verb|MAX_DIVISORS| (\csource{find_points.c}{864}{880}, the test at
\csource{find_points.c}{881}{881}, which sets the flag \verb|RATPOINTS_USE_SQUARES1| at
\csource{find_points.c}{885}{885}). Then the denominators are $e k^2 \le H$, divisor $e$ by divisor
in the order of binary counting over the primes of $c_d$, ascending, and $k$ ascending for each
(\csource{find_points.c}{3937}{3940}). Stage~E proved this order through the definitions of the
model (Theorem~\ref{thm:findPoints_order}); G states it in terms of the polynomial.

\paragraph{Why the list does not change.} With a callback that appends each point to an array, with
the weight $1$, and never sets \verb|quit|, the delivery of a list is that list and the count its
length; so \verb|findPoints|, the search with that callback, returns the list of before G when
neither option of the output is set. For any callback the search delivers the list of
\verb|findPoints| for the same options, up to and including the first call that stops it: the stops
only cut that list. The options of the output act on the list itself, for \verb|findPoints| as well,
by the filter of \verb|RATPOINTS_NO_Y| and, under \verb|RATPOINTS_NO_CHECK|, by the survivors in
place of the points. So the theorems on the points are stated for \verb|findPoints|, and the
protocol theorem (Theorem~\ref{thm:protocol}) carries them over to the search with any callback.

\section{The delivery}\label{sec:delivery}

\begin{definition}[The callback and the delivery]\label{def:delivery}
  \lean{Ratpoints.Answer, Ratpoints.Delivery, Ratpoints.Delivery.call, Ratpoints.CheckFlags, Ratpoints.Delivery.callMapped, Ratpoints.appendPoint}
  \leanok
  What one call of the callback gives: the weight that it returns, an integer with no bound, whether
  it set \verb|*quit|, and the data after the call (\csource{ratpoints.h}{436}{437},
  \csource{ratpoints.h}{441}{442}, the prototypes of the callback, which
  Definition~\ref{def:findPointsWith} shares); the callback is a function of $x$, $z$, $y$ and the
  data, in the order of the C code, that returns what the call gives, with no argument for
  \verb|quit|, which is clear at every call. The delivery of a search: the data of the callback, the
  count, the sum of the weights so far, and the flag of the stop; one call
  \verb|total += process(x, z, y, info, &quit)| on a delivery adds the weight to the count and takes
  the flag and the data that the call leaves. The bits of the flags that the exact check and the
  points at infinity read, \verb|RATPOINTS_REVERSED|, \verb|RATPOINTS_NO_CHECK| and
  \verb|RATPOINTS_NO_Y|, as three Booleans, all clear by default (\csource{ratpoints.h}{414}{414},
  which Definition~\ref{def:Outcome} shares). One call with the point $(a : y : b)$ of the
  polynomial searched, mapped back to the caller's curve when the flag of the reversal is set: with
  $(b : y : a)$ for $a \ge 0$ and $(-b : y : -a)$ for $a < 0$, as $\mathrm{swapPoint}$ maps the
  point (Definition~\ref{def:swapPoint}), the three branches that the exact check writes at each of
  its calls (lines 5312-5316, 5321-5325 and 5331-5335) as one definition. And the callback of the
  model that appends each point $(x, y, z)$ to an array, with the weight $1$, and never stops the
  search. The delivery, the call and the mapped call have no C ranges of their own: the lines they
  stand for (lines 4570-4571 and 3129 for the count and the two flags of the stop, \verb|quit| and
  \verb|wk->quit|, and the branches of the map) belong to their callers; the appending callback
  models no C code.
\end{definition}

\begin{definition}[The reference of the delivery]\label{def:deliver}
  \uses{def:delivery}
  \lean{Ratpoints.Delivery.deliver}
  \leanok
  The delivery of the points of a list $L$, in their order, to the callback on a delivery $r$: for
  each point $(x, y, z)$ of $L$, while the flag of $r$ is clear, the call with $(x : y : z)$; the
  points after a call that sets the flag, and all of them when it is set at the start, are passed
  over. It is a reference of the model, as the brute-force reference is, and models no C code: the
  search delivers each point where it finds it and tests the flag where the C code tests
  \verb|quit|, and this fold over the list of the points of \verb|findPoints| is what that delivery
  is compared with (Theorem~\ref{thm:protocol}).
\end{definition}

\begin{lemma}[The delivery of a list]\label{lem:deliver}
  \uses{def:deliver}
  \lean{Ratpoints.Delivery.deliver_nil, Ratpoints.Delivery.deliver_cons, Ratpoints.Delivery.deliver_append, Ratpoints.Delivery.deliver_flatMap, Ratpoints.Delivery.deliver_cons_of_quit_eq_false, Ratpoints.Delivery.deliver_singleton_of_quit_eq_false, Ratpoints.Delivery.deliver_of_quit, Ratpoints.Delivery.deliver_append_of_quit, Ratpoints.Delivery.deliver_eq_of_prefix, Ratpoints.Delivery.deliver_flatMap_eq_of_prefix, Ratpoints.Delivery.ite_quit_eq_deliver_append, Ratpoints.Delivery.deliver_appendPoint, Ratpoints.Delivery.deliver_eq_foldl_call, Ratpoints.Delivery.exists_prefix_deliver_eq, Ratpoints.Delivery.total_deliver_eq_sum, Ratpoints.Delivery.exists_prefix_total_deliver_eq_sum}
  \leanok
  The delivery of the empty list leaves a delivery as it is, and that of $p :: L$ delivers $L$ after
  the call with $p$, when the flag is clear; the delivery of $L_1 \mathbin{+\!\!+} L_2$ delivers
  $L_2$ after $L_1$, and that of a \verb|flatMap| is the fold of the deliveries of its parts. A
  delivery whose flag is set stays as it is, so that nothing changes after a list whose delivery
  sets the flag: for a prefix $M$ of $L$ whose delivery sets it, the delivery of $L$ is that of $M$,
  also for the prefix of the index list of a \verb|flatMap|; and a test of the flag after the
  delivery of $L_1$, followed by the delivery of $L_2$ when it is clear, is the delivery of
  $L_1 \mathbin{+\!\!+} L_2$. With the appending callback, from a delivery whose flag is clear, the
  array grows by $L$ and the count by its length. For every list $L$ there is a prefix $M$ of $L$
  whose delivery is that of $L$, such that the flag is clear before each point of $M$, so that each
  of them is called, and $M = L$ unless the delivery of $M$ sets the flag; the count grows by the
  sum of the weights of these calls, each taken on the data that the calls before it left.
\end{lemma}

\begin{proof}
  \leanok
  By induction on the list, from the fold. The count as a sum over the indices of $M$ needs
  Mathlib's sums over \verb|Fin|, in a file of its own, and the rest is core.
\end{proof}

\section{The exact check and the points at infinity}\label{sec:checkDelivery}

The exact check (Definition~\ref{def:accepted}) takes the callback, the flags and a delivery: unless
$\mathrm{noCheck}$, for $v$ the value of Horner's rule in $a$ on the products, times $b$ for odd
$d$, and $y = \lfloor \sqrt{v} \rfloor$, when $v \ge 0$ and $v = y^2$, the mapped call with $(a : y
: b)$, and then, unless that call stopped the search, $y = 0$ or $\mathrm{noY}$, the mapped call
with $(a : -y : b)$; with $\mathrm{noCheck}$ one mapped call with $(a : 0 : b)$, whether or not it
gives a point. The points at infinity (Definition~\ref{def:pointsAtInfinity}) take them too: when
there are points at infinity, for odd $d$ or $c_d$ a square, with the candidate $(1, 0)$, or $(0,
1)$ when the flag of the reversal is set, one call with $y = 0$ under $\mathrm{noCheck}$ or for odd
$d$, and for even $d$ the call with $s = \lfloor \sqrt{c_d} \rfloor$ and then, unless that call
stopped the search or $\mathrm{noY}$, the call with $-s$, with no test $s \ne 0$, as the C code has
none: the case $c_d = 0$ for even $d$, which would deliver $(1, 0, 0)$ twice, cannot arise after the
refusals, in the C code as in the model, since the leading coefficient of the polynomial searched is
not zero. The return after a stop at infinity is made by the search. The contracts of the two say
what they deliver to a callback that never stops the search, a list that depends on the flags; the
sieve passes these lists on, and they replace the check of a numerator with the flag of the reversal
(Definition~\ref{def:checkNum}) in its contracts.

\begin{definition}[The points of a numerator with the flags]\label{def:checkNumWith}
  \uses{def:delivery, def:bruteForce, def:swapPoint, def:reversal}
  Proposed names: \lname{Ratpoints.Sieve.checkPointList}, \lname{Ratpoints.Sieve.checkNumWith},
  \lname{Ratpoints.Sieve.pointsAtInfinityList}. The points of the exact check of the numerator $a$
  over the denominator $b$, for the products $w$, the factor $f$ and the flags: under
  $\mathrm{noCheck}$ the survivor $(a, 0, b)$; otherwise $(a, y, b)$ and, for $y \ne 0$, $(a, -y,
  b)$ when the value of Horner's rule in $a$ on $w$, times $f$, is $y^2$ with $y \ge 0$, the point
  with $y < 0$ left out under $\mathrm{noY}$; each mapped by $\mathrm{swapPoint}$ when the flag of
  the reversal is set. The check of a numerator with the flags: these points when $\gcd(a, b) = 1$,
  and none otherwise, the reference of the test of a candidate and the exact check with the flags.
  The points at infinity with the flags, for the coefficient list $c$ of degree $d$: when $d$ is odd
  or $c_d$ is a square, $(1, 0, 0)$ under $\mathrm{noCheck}$ or for odd $d$, and otherwise $(1, s,
  0)$ for $s = \lfloor \sqrt{c_d} \rfloor$ followed, unless $\mathrm{noY}$, by $(1, -s, 0)$, each
  mapped by $\mathrm{swapPoint}$ when the flag of the reversal is set, which gives $(0, y, 1)$; none
  when $d$ is even and $c_d$ is not a square, the square test that of the reversal check
  (Definition~\ref{def:reversal}).
\end{definition}

\begin{lemma}[The flags of the check]\label{lem:checkNumWith}
  \uses{def:checkNumWith, def:checkNum}
  Proposed names: \lname{Ratpoints.Sieve.checkNumWith\_of\_noCheck\_eq\_false},
  \lname{Ratpoints.Sieve.checkNumWith\_of\_noCheck\_eq\_true},
  \lname{Ratpoints.Sieve.checkNumWith\_eq\_map},
  \lname{Ratpoints.Sieve.mem\_checkNumWith\_of\_noCheck},
  \lname{Ratpoints.Sieve.map\_checkNumWith\_sublist},
  \lname{Ratpoints.Sieve.flatMap\_checkNumWith\_of\_noCheck\_eq\_false},
  \lname{Ratpoints.Sieve.pointsAtInfinityList\_of\_noCheck\_eq\_false},
  \lname{Ratpoints.Sieve.pointsAtInfinityList\_of\_noCheck\_eq\_true},
  \lname{Ratpoints.Sieve.pointsAtInfinityList\_eq\_map},
  \lname{Ratpoints.Sieve.eq\_of\_mem\_pointsAtInfinityList\_of\_noCheck},
  \lname{Ratpoints.Sieve.map\_pointsAtInfinityList}, \lname{Ratpoints.Sieve.fst\_snd\_swapPoint},
  \lname{Ratpoints.Sieve.filter\_map\_swapPoint}, \lname{Ratpoints.Sieve.swapPoint\_fst\_zero\_snd}.
  Without $\mathrm{noCheck}$ the check of a numerator with the flags is the check of the reference
  with the flag of the reversal, without the points with $y < 0$ under $\mathrm{noY}$; so over a
  list of numerators it gives the points that the brute-force reference finds over $b$, mapped after
  a reversal and filtered, the map keeping $y$, so that it commutes with the filter. Under
  $\mathrm{noCheck}$ it is the mapped survivor $(a, 0, b)$ when $\gcd(a, b) = 1$ and nothing
  otherwise, and a triple lies in it exactly then. With the flag of the reversal set it gives the
  points with it clear, mapped. The points with the exact check on and $\mathrm{noY}$ set, with $y$
  replaced by $0$, are a sublist of those under $\mathrm{noCheck}$ with the same flag of the
  reversal, for either setting of $\mathrm{noY}$ there. The points at infinity with the flags:
  without $\mathrm{noCheck}$, for odd $d$ or $c_d \ne 0$, the points of the candidate $(1, 0)$ of
  the reference, mapped and filtered; under $\mathrm{noCheck}$ the one mapped point $(1, 0, 0)$ when
  there are points at infinity, every member being that point; with the flag of the reversal set,
  those with it clear, mapped; and those with the exact check on and $\mathrm{noY}$ set, with $y$
  replaced by $0$, are those under $\mathrm{noCheck}$.
\end{lemma}

\begin{proof}
  By the cases of the flags. The points of the reference over one numerator are $(a, y, b)$ and $(a,
  -y, b)$, or $(a, 0, b)$, or none, so that the point with $y \ge 0$, with $y$ replaced by $0$, is
  the survivor; the map of the reversal keeps $y$ and maps $(a, 0, b)$ to the image of the survivor.
\end{proof}

\begin{lemma}[The check and the points at infinity deliver their lists]\label{lem:checkDelivery}
  \uses{def:checkNumWith, lem:deliver, def:accepted, def:pointsAtInfinity}
  Proposed names: \lname{Ratpoints.Sieve.callMapped\_eq\_deliver},
  \lname{Ratpoints.Sieve.checkPoint\_eq\_deliver},
  \lname{Ratpoints.Sieve.checkPoint\_eq\_deliver\_acceptedPoints},
  \lname{Ratpoints.Sieve.deliver\_acceptedPoints\_of\_not\_accepted},
  \lname{Ratpoints.Sieve.pointsAtInfinity\_eq\_deliver}. On a delivery whose flag is clear: the
  mapped call is the delivery of the one mapped point; the exact check delivers its points
  (\verb|checkPointList|), the call with $-y$ after a call that stops the search being the delivery
  that passes over the rest; a numerator that passes the test of a candidate delivers the points of
  the extraction, one that does not delivers nothing; and the points at infinity deliver their list
  with the flags.
\end{lemma}

\begin{proof}
  Each call is the delivery of a singleton (Lemma~\ref{lem:deliver}), and the test of the flag
  between the two calls of a pair is the delivery of the concatenation.
\end{proof}

\section{The sieve on the delivery}\label{sec:sieveDelivery}

Every component of the sieve from the walk through the set bits of a word up to the run over the
denominators takes the callback, the flags and a delivery, in place of the flag of the reversal and
of the array of the points or the list that it returned, hands them on to the components it calls,
and returns after a call that stops the search, as the C code returns at its tests of \verb|quit|.
The walk (Definition~\ref{def:walkWord}) returns after a check at which the callback set the flag,
as \verb|_ratpoints_sift0| does at its three tests, lines it already owns; the two arms of a bit
array through the second stage and the extraction (Definition~\ref{def:siftBitArray}) test the flag
after each walk and the scan (Definition~\ref{def:scanPiece}) after each bit array, where the C code
returns from \verb|_ratpoints_sift0| from within the walk, which makes no call that the C code does
not make and skips none; one piece (Definition~\ref{def:siftPiece}) returns what the scan returns;
the bits of one denominator (Definition~\ref{def:siftBits}) test it after each piece, as \verb|sift|
does (line 3611), one denominator (Definition~\ref{def:sift}) returns the delivery with the store of
the tables of rows, and the run (Definition~\ref{def:sieveRun}) tests it after each denominator, for
the returns of \verb|run_block| in each shape and the end of the loop over the blocks, lines it
already owns. The stop test after each walk, three of them, would take \verb|siftBitArray| from $18$
to $21$ lines, past the bound of $20$ lines on a definition, so the two arms of a bit array through
the second stage and the extraction become functions of their own,
\lname{Ratpoints.Sieve.siftBitArrayByWords} for the arm of the builds with
\verb|USE_LONG_IN_PHASE_2| (\csource{sift.c}{1126}{1134}, \csource{sift.c}{1137}{1139},
\csource{sift.c}{1154}{1159}, \csource{sift.c}{1172}{1173}) and
\lname{Ratpoints.Sieve.siftBitArrayWhole} for that of the default build
(\csource{sift.c}{1204}{1210}, \csource{sift.c}{1213}{1226}, \csource{sift.c}{1257}{1264},
\csource{sift.c}{1293}{1296}, with the lines of \verb|EXT0| and \verb|EXT| of \verb|rp-private.h|),
each keeping the first numerator of the bit array and the step from one word to the next where the C
code computes them, as \verb|siftBitArray| does today; \verb|siftBitArray| chooses between them by
its argument and keeps \csource{sift.c}{1094}{1111}, and its deviation gains ``each arm a function
of its own''. Every C range of the three is one that \verb|siftBitArray| holds today.

The contracts keep their hypotheses, gain the precondition that the flag of the delivery is clear,
and conclude that the delivery after the component is the delivery of a list
(Definition~\ref{def:deliver}) from the delivery before it: the list that the contract gave before
G, with the points of the extraction with the flags (\verb|acceptedPoints|,
Definition~\ref{def:accepted}, the points of the check with the flags of
Definition~\ref{def:checkNumWith} for a numerator that passes the test of a candidate) in place of
those with the flag of the reversal, from the walk (\lname{Ratpoints.Sieve.walkWord\_spec}), the
arms (\lname{Ratpoints.Sieve.siftBitArrayByWords\_spec},
\lname{Ratpoints.Sieve.siftBitArrayWhole\_spec}, \lname{Ratpoints.Sieve.siftBitArray\_spec}), the
scan (\lname{Ratpoints.Sieve.scanPiece\_spec}) and one piece
(\lname{Ratpoints.Sieve.siftPiece\_spec}) up to the bits of one denominator
(\lname{Ratpoints.Sieve.siftBits\_spec}). For one denominator (\lname{Ratpoints.Sieve.sift\_spec})
the list is that of the numerators of $[-H, H]$ that pass the class bit of the word $w$ that the
class packs and survive the entries of all three stages, each with the points of its check with the
flags (\lname{Ratpoints.Sieve.flatMap\_acceptedPoints\_eq\_flatMap\_checkNumWith}), and the store
after the denominator fits the entries again; for the run (\lname{Ratpoints.Sieve.sieveRun\_spec},
with its equation \lname{Ratpoints.Sieve.run\_sieveRun} and the form
\lname{Ratpoints.Sieve.run\_sieveRun\_of\_eq\_filter} for a list of the shape that is a list $l$
filtered by a test, which delivers over $l$) it is the concatenation of these lists over the
denominators of the shape that pass its test. Without $\mathrm{noCheck}$ the words still compute the
Booleans (Lemma~\ref{lem:sift_eq_checkDen},
\lname{Ratpoints.Sieve.run\_sift\_eq\_deliver\_checkDen}): one denominator delivers the points of
the check of stage~C1 on the sieve of stage~C1 whose numerators pass the class bit, mapped after a
reversal and filtered by $\mathrm{noY}$.

The loops that may return early carry invariants of their own, each in the node of the contract that
uses it. Each loop of the sieve that delivers points, over the set bits of a word, the words of a
bit array, the bit arrays of a piece, the pieces of an interval and the denominators of a run, keeps
an invariant with two cases: the flag is clear and the delivery is that of the points of the steps
done so far, or the loop has returned at a stop, with the flag set, and the delivery is that of the
points of the steps up to the one that stopped it, which is that of the points of all the steps
(Lemma~\ref{lem:deliver} on a prefix). These are, for the walk,
\lname{Ratpoints.Sieve.WalkDeliveryInvariant} and \lname{Ratpoints.Sieve.WalkDeliveryWhileInvariant}
(Definition~\ref{def:walkWord}); for the words of a bit array,
\lname{Ratpoints.Sieve.SiftBitArrayLoopInvariant} and
\lname{Ratpoints.Sieve.SiftBitArrayWalkLoopInvariant}, with
\lname{Ratpoints.Sieve.SiftBitArrayInvariant.eq\_deliver\_of\_quit},
\lname{Ratpoints.Sieve.SiftBitArrayWalkInvariant.eq\_deliver\_of\_quit} and
\lname{Ratpoints.Sieve.eq\_deliver\_of\_quit\_of\_forall\_eq} for the stop
(Lemma~\ref{lem:siftBitArray}); for the bit arrays of a piece,
\lname{Ratpoints.Sieve.ScanPieceInvariant.eq\_deliver\_of\_quit} (Definition~\ref{def:scanPiece});
for the masked bits of one piece, \lname{Ratpoints.Sieve.IsMasked.eq\_deliver}
(Definition~\ref{def:siftPiece}); for the pieces of an interval,
\lname{Ratpoints.Sieve.SiftInvariant.eq\_deliver\_of\_quit} (Lemma~\ref{lem:siftBits}); and for the
denominators of a run, \lname{Ratpoints.Sieve.SieveRunLoopInvariant} and
\lname{Ratpoints.Sieve.SieveRunInvariant.eq\_deliver\_of\_quit} (Definition~\ref{def:sieveRun}).
What they share is the following lemma.

\begin{lemma}[The steps done come first]\label{lem:sieveStops}
  \uses{lem:numClasses, def:andRows, def:rbaPack}
  Proposed names: \lname{Ratpoints.Sieve.flatMap\_numerator\_range\_prefix},
  \lname{Ratpoints.Sieve.flatMap\_range\_rbaPack\_mul\_prefix},
  \lname{Ratpoints.Sieve.flatMap\_filter\_range\_prefix}, \lname{List.range\_prefix\_range}. The
  points of the first steps of a loop of the sieve are a prefix of the points of all its steps: the
  points of the numerators of the set bits of the first $k \le m$ words of a bit array among those
  of the first $m$, those of the bits of the first $\mathrm{pos} \le n$ bit arrays of a piece among
  those of the first $n$, those of the bits up to the word number $w_0$ among those up to the word
  number $w_0' \ge w_0$, and the numbers below $k \le m$ among those below $m$. So in both cases of
  an invariant the delivery at the end is that of the points of all the steps, and each contract
  concludes with the delivery of its whole list.
\end{lemma}

\begin{proof}
  The lists of the first steps are the first parts of a \verb|flatMap| over a range, and a range
  below $k \le m$ is the first part of the range below $m$. The invariants are built, with the
  program logic of Lean core, with the case of the early return that core provides for the
  \verb|for| loops, and with a case for the exit for the \verb|while| and \verb|repeat| loops; their
  verification conditions are closed by lemmas stated in their shape, with the invariants
  irreducible after their lemmas, and the case of a stop by Lemma~\ref{lem:deliver} on the prefix.
\end{proof}

\section{The search on the delivery}\label{sec:searchDelivery}

The search of the polynomial handed on (Definition~\ref{def:searchPoints}) takes the callback, the
flags and a delivery and runs in the order of \verb|find_points_work_1|: the flags of the tests on
the denominators and the divisors of the leading coefficient, the 2-adic information with its exit,
what \verb|sieving_info| computes with its exit, the switch of the checked loop, then the points at
infinity, the return when the callback has stopped the search (\csource{find_points.c}{5100}{5106},
which the search takes over from the points at infinity, with the return when the Sturm step leaves
nothing to search, not met before stage~H), and the sieving, a function of its own. Both exits
return the delivery as it came. The search returns what a search gives back: the data and the count
of the delivery with the bits of the flags written back.

\begin{definition}[The sieving]\label{def:sieving}
  \uses{def:sieveRun, def:classOffsets, def:Shape, def:delivery, def:arraySize}
  Proposed name: \lname{Ratpoints.sieving}. The block that the C code introduces as ``now do the
  sieving'': an array of rows of offsets, one for each packing, of one offset for each sieve entry,
  and the row offsets of the classes in it, with the inverses of the strides in the entries
  (Definition~\ref{def:classOffsets}); the shape that \verb|run_setup| chooses
  (Definition~\ref{def:Shape}); and the run over the denominators $1, \dots, H$ of the shape with
  the cut points of the parameters, pieces of at most $\mathrm{arraySizeAt}(H)$ bit arrays, the
  flags and the delivery, which returns when the callback stops the search
  (\csource{find_points.c}{4573}{4573}, \csource{find_points.c}{5116}{5116},
  \csource{find_points.c}{5119}{5119}, \csource{find_points.c}{5132}{5132},
  \csource{find_points.c}{5137}{5139}, \csource{find_points.c}{5167}{5167}, which move from the
  search, Definition~\ref{def:searchPoints}).
\end{definition}

\begin{definition}[What a search gives back]\label{def:Outcome}
  Proposed name: \lname{Ratpoints.Outcome}. The data of the callback after the search, the count,
  which \verb|find_points| returns, and the three bits of \verb|args->flags| that the search writes
  back and the program reads after it: \verb|RATPOINTS_REVERSED|, \verb|RATPOINTS_USE_SQUARES|, set
  for odd degree and $|c_d| = 1$, and \verb|RATPOINTS_USE_SQUARES1|, set for odd degree otherwise
  when the divisor set-up succeeds (\csource{ratpoints.h}{377}{377},
  \csource{ratpoints.h}{414}{414}, \csource{ratpoints.h}{416}{417}). The search sets the last two
  before its exits and before the points at infinity, so that they are written back on every path
  that does not refuse the input. Nothing else is written back: not the reversed coefficients and
  the degree, which nothing reads after the search; not the other internal bits of the flags,
  \verb|RATPOINTS_CHECK_DENOM|, \verb|RATPOINTS_USE_JACOBI| and \verb|RATPOINTS_NO_REVERSE_AUTO|,
  which the program does not read either; not the search intervals, \verb|args->num_inter| and
  \verb|args->domain|, which the search sets to $[-H, H]$ when none is given and the Sturm step
  narrows, and which message 8 reads after the search (stage~H); and not the numbers of the moduli
  used (stage~I).
\end{definition}

\begin{definition}[The search with a callback]\label{def:findPointsWith}
  \uses{def:Outcome, def:delivery, def:searchPoints, def:Options, def:reversal, def:sturmChain,
    def:useCLong, def:maxDegree, def:Params}
  Proposed name: \lname{Ratpoints.findPointsWith}. The library function \verb|find_points| with the
  callback and its data, as \verb|find_points|, \verb|find_points_work| and
  \verb|find_points_work_1| do it: the refusals of the search for the points in the order of the
  program, more than $\mathrm{maxDegree} + 1$ coefficients, fewer than two, $H = 0$ or $H \ge
  2^{63}$, and a form whose leading zeros lower $\lceil n/2 \rceil$; the reversal check with the
  switch of the options (Definition~\ref{def:reversal}), which may refuse the form or reverse the
  polynomial; with the Sturm step on, the refusal when the remainder chain of the polynomial
  searched is $\mathrm{none}$ (Definition~\ref{def:sturmChain}); and otherwise the search of the
  polynomial handed on, with the flag of the reversal and the options of the output of the options
  as its flags, from the data of the caller, the count $0$ and the flag of the stop clear, which
  returns what the search gives back (\csource{find_points.c}{4523}{4528},
  \csource{find_points.c}{4546}{4546}, \csource{find_points.c}{4557}{4559},
  \csource{find_points.c}{4564}{4565}, \csource{find_points.c}{4567}{4571},
  \csource{find_points.c}{4578}{4585}, \csource{find_points.c}{4606}{4606},
  \csource{find_points.c}{4624}{4633}, \csource{find_points.c}{4658}{4658},
  \csource{find_points.c}{4857}{4892}, \csource{find_points.c}{5219}{5220},
  \csource{find_points.c}{5226}{5229}, \csource{find_points.c}{5235}{5235},
  \csource{find_points.c}{5241}{5242}, \csource{ratpoints.h}{362}{362}, the links of the search for
  the points until now, and \csource{ratpoints.h}{427}{430}, \csource{ratpoints.h}{436}{437},
  \csource{ratpoints.h}{441}{442}, the codes of the refusals and the prototypes, and
  \csource{ratpoints.h}{409}{411}, the mask of the bits of the input that line 4606 keeps). The
  search for the points (Definition~\ref{def:findPoints}) becomes its instance with the appending
  callback from the empty array, the list of the array that it returns, or its refusal; it models no
  C code and holds no C range.
\end{definition}

The theorems of the search speak about the list of the points that it delivers for a test of the
sieve: the sieve of the search lets a numerator through by its class word, the moduli of three
stages and the test for common factors, for a given polynomial and parameters, and what a contract
can say of the test is that it lets the numerator of every point through. One test serves every
callback and every setting of the flags, since the components that decide it, the classes, the word
\verb|den_bits|, the exits and what \verb|sieving_info| returns, depend on neither; this is what
relates the output under \verb|RATPOINTS_NO_CHECK| to the points.

\begin{definition}[The points that the search delivers]\label{def:searchList}
  \uses{def:checkNumWith, def:searchDens, def:bruteForce}
  Proposed names: \lname{Ratpoints.searchList}, \lname{Ratpoints.IsSoundTest}. For the coefficient
  list $c$ of its degree $d$, the height $H$, a test $S(b, a)$ of the numerators $a$ over the
  denominators $b$ and the flags: the points at infinity with the flags, then, over the denominators
  $b$ of the search in the order of its loop (Definition~\ref{def:searchDens}), the points of the
  check with the flags of the numerators $a$ of $-H, \dots, H$ that pass $S(b, \cdot)$, with the
  products of $c$, $d$ and $b$ and the factor $b$ for odd $d$ and $1$ for even $d$. A test is
  \emph{sound} when, over each denominator $b$ of the search, it lets through every numerator $a$ of
  $-H, \dots, H$ for which the brute-force reference finds a point $(a, y, b)$.
\end{definition}

\begin{lemma}[The search delivers the points of a sound test]\label{lem:searchList}
  \uses{def:searchList, def:sieving, def:searchPoints, def:Outcome, lem:sieveStops,
    thm:run_reference, lem:soundRun, lem:checkDelivery, lem:checkNumWith, lem:shape_dens,
    lem:twoAdic_exit, lem:sievingInfo, lem:examinePrime_none}
  Proposed names: \lname{Ratpoints.exists\_isSoundTest\_run\_searchPoints},
  \lname{Ratpoints.exists\_isSoundTest\_run\_sieveRun},
  \lname{Ratpoints.run\_searchPoints\_eq\_deliver}, \lname{Ratpoints.run\_searchPoints\_flags},
  \lname{Ratpoints.searchList\_of\_noCheck\_eq\_false}, \lname{Ratpoints.searchList\_false},
  \lname{Ratpoints.isSoundTest\_false\_of\_forall\_not\_isPoint},
  \lname{Ratpoints.pointsAtInfinityList\_eq\_nil\_of\_forall\_not\_isPoint},
  \lname{Ratpoints.Sieve.pointsAtInfinityList\_eq\_filter\_map\_pointsOfDen},
  \lname{Ratpoints.Sieve.pointsAtInfinityList\_eq\_nil\_of\_pointsOfDen\_eq\_nil},
  \lname{Ratpoints.BruteForce.pointsOfDen\_zero\_singleton\_one},
  \lname{Ratpoints.BruteForce.pointsOfDen\_ne\_nil\_of\_mem}. For well-formed parameters and $1 \le
  H < 2^{63}$ there is a sound test $S$ such that the search of $c$ of its degree, with any
  callback, any flags and a delivery whose flag is clear, gives the data and the count of the
  delivery of the list of the points of Definition~\ref{def:searchList} for $S$ and these flags; so
  it gives those of the delivery of the list that the search with the appending callback gives, from
  the empty array, for the same flags. The run of the shape delivers the part over the denominators
  for a sound test, whatever the callback and the flags. The bits of the flags that the search
  writes back are the flag of the reversal as it came, the squares for odd degree and $|c_d| = 1$,
  and the success of the divisor set-up for odd degree and $|c_d| \ne 1$, whatever the callback.
  Without $\mathrm{noCheck}$, for a sound test and $c$ of positive degree, the points delivered are
  those that the brute-force reference finds over the candidate $(1, 0)$ and over the denominators
  of the search, in this order, mapped by $\mathrm{swapPoint}$ after a reversal, without those with
  $y < 0$ under $\mathrm{noY}$. When the curve has no points, the test that lets nothing through is
  sound and there are no points at infinity.
\end{lemma}

\begin{proof}
  \uses{lem:flatMap_sound, lem:survives, thm:bruteForce}
  A lemma over the paths of the search gives its data and its count along each path. At the 2-adic
  exit and the exit of \verb|sieving_info| the curve has no points (Lemmas~\ref{lem:twoAdic_exit}
  and~\ref{lem:examinePrime_none}), so that the test that lets nothing through is sound and nothing
  is delivered; otherwise the points at infinity come first (Lemma~\ref{lem:checkDelivery}), the
  search returns after them at a stop, and the run delivers the points of the check with the flags
  of the numerators that pass the test of the shape, the class words and the entries (the contract
  of the run, Definition~\ref{def:sieveRun}, and Theorem~\ref{thm:run_reference}), which let every
  point through, by what the search hands to the run (Lemma~\ref{lem:soundRun}) and by
  Lemma~\ref{lem:shape_dens}. The case analysis is on the values of the components that decide the
  test, which depend neither on the callback nor on the flags, so that one test serves them all.
\end{proof}

\begin{theorem}[The protocol of the delivery]\label{thm:protocol}
  \uses{def:findPointsWith, def:findPoints, lem:searchList, lem:deliver, lem:reversal_swap}
  Proposed names: \lname{Ratpoints.map\_findPointsWith\_eq\_map\_findPoints},
  \lname{Ratpoints.findPointsWith\_eq\_ok\_iff}, \lname{Ratpoints.findPointsWith\_eq\_error\_iff},
  \lname{Ratpoints.exists\_findPoints\_eq\_error\_iff},
  \lname{Ratpoints.map\_findPointsWith\_flags}. For well-formed parameters, every option and every
  callback with its data: $\mathrm{findPointsWith}$ refuses exactly the inputs that
  $\mathrm{findPoints}$ refuses, with the same errors, and otherwise gives the data and the count of
  the delivery (Definition~\ref{def:deliver}) of the list of $\mathrm{findPoints}(c, H, P, o)$ to
  the callback, from its data, the count $0$ and the flag of the stop clear: the calls of the
  callback are those of the points of that list, in its order, up to and including the first that
  stops the search, and the count is the sum of their weights (Lemma~\ref{lem:deliver}). What
  $\mathrm{findPointsWith}$ returns is the outcome of the search of the polynomial that the reversal
  check hands on, once none of the refusals applies; its refusals and the bits of the flags that it
  writes back do not depend on the callback or its data.
\end{theorem}

\begin{proof}
  The refusals are the same tests on both sides; after them both searches run on the polynomial
  handed on, with its degree, which the refusals leave positive, and Lemma~\ref{lem:searchList}
  gives the data and the count. The theorem needs the well-formedness of the parameters, since the
  equality comes from the contracts of the search, which need it, and not from an argument that the
  control flow depends on the callback only through the flag.
\end{proof}

\begin{lemma}[The list of the search]\label{lem:findPoints_searchList}
  \uses{lem:searchList, thm:protocol, def:Options, lem:reversal_swap}
  Proposed name: \lname{Ratpoints.exists\_searchList\_of\_findPoints\_eq\_ok}. For well-formed
  parameters, if $\mathrm{findPoints}(c, H, P, o)$ returns the list $L$, then the reversal check
  hands on a polynomial $c'$ of positive degree with the flag $\mathrm{rev}$ of the reversal, $1 \le
  H$, and there is a sound test $S$ such that $L$ is the list of the points that the search of $c'$
  delivers for $S$ (Definition~\ref{def:searchList}) with the flags $\mathrm{rev}$,
  $\mathrm{noCheck}$ and $\mathrm{noY}$ of $o$, and such that, for every choice of the options of
  the output with the other options unchanged, $\mathrm{findPoints}$ returns the list for $S$ with
  those flags.
\end{lemma}

\begin{proof}
  Lemma~\ref{lem:searchList} for the polynomial handed on, with the appending callback; the test
  does not depend on the flags.
\end{proof}

\section{The options of the output}\label{sec:outputOptions}

The options of the search (Definition~\ref{def:Options}) gain the fields $\mathrm{noCheck}$ and
$\mathrm{noY}$, false by default, for the bits \verb|RATPOINTS_NO_CHECK| and \verb|RATPOINTS_NO_Y|
of the flags (\csource{ratpoints.h}{403}{405}, with \verb|RATPOINTS_NO_REVERSE|, which the switch of
the reversal already stands for), which the search with a callback copies into the flags of the
check. With them the list of $\mathrm{findPoints}$ is, without $\mathrm{noCheck}$, the list of the
points with those of $y < 0$ left out under $\mathrm{noY}$, and under $\mathrm{noCheck}$ the list of
the survivors with $y = 0$, which depends on the parameters. The main theorems
(Theorem~\ref{thm:findPoints}) take them into account: for $\mathrm{noCheck}$ false, a triple lies
in the list exactly when it is a normalized point of height at most $H$, with $y \ge 0$ under
$\mathrm{noY}$ (\lname{Ratpoints.mem\_iff\_isPoint\_of\_findPoints\_eq\_ok}); the list has no
duplicates for every option (\lname{Ratpoints.nodup\_of\_findPoints\_eq\_ok}); and its set of
elements is $\mathrm{points}(c, H)$ when both options of the output are off
(\lname{Ratpoints.toFinset\_eq\_points\_of\_findPoints\_eq\_ok}). The theorem of the output of the
search (Theorem~\ref{thm:findPoints_order}, \lname{Ratpoints.exists\_eq\_of\_findPoints\_eq\_ok})
gives, for $\mathrm{noCheck}$ false, the reference of the polynomial searched filtered by
$\mathrm{noY}$, and the reference itself and the permutation of it
(Theorem~\ref{thm:findPoints_perm}, \lname{Ratpoints.perm\_bruteForce\_of\_findPoints\_eq\_ok}) take
both options of the output off. The refusal theorems (Theorem~\ref{thm:refusals}) keep their
statements. Under $\mathrm{noCheck}$ the following lemma and theorem say what the list holds.

\begin{lemma}[The survivors under NO\_CHECK]\label{lem:noCheckList}
  \uses{def:searchList, lem:checkNumWith, lem:swapPoint_bijection, def:searchDens, lem:shape_dens}
  Proposed names: \lname{Ratpoints.exists\_eq\_of\_mem\_searchList\_of\_noCheck},
  \lname{Ratpoints.nodup\_searchList\_of\_noCheck}, \lname{Ratpoints.map\_searchList\_sublist},
  \lname{Ratpoints.searchList\_eq\_map}, \lname{Ratpoints.nodup\_searchDens},
  \lname{List.Nodup.flatMap}. Under $\mathrm{noCheck}$ and for $1 \le H$, every member of the list
  that the search of $c$ delivers for any test is $(x, 0, z)$ for a pair $(x, z)$ normalized of
  height at most $H$ (Lemma~\ref{lem:swapPoint_bijection}), and the list has no duplicates. For any
  test, with the flag of the reversal set, the list is that with it clear, mapped; and the list with
  the exact check on and $\mathrm{noY}$ set, with $y$ replaced by $0$, is a sublist of the list
  under $\mathrm{noCheck}$ with the same flag of the reversal. The denominators of the search have
  no duplicates.
\end{lemma}

\begin{proof}
  The denominators of the search lie in $1, \dots, H$, as the search uses already. The survivors of
  distinct denominators, or of distinct numerators of one denominator, are distinct pairs; a list of
  lists without duplicates whose members carry the index of their part has no duplicates
  (\verb|List.Nodup.flatMap|, a candidate for Mathlib); and the map of the reversal is injective on
  normalized pairs (Lemma~\ref{lem:swapPoint_bijection}). The sublist holds per numerator and at
  infinity (Lemma~\ref{lem:checkNumWith}), and both lists run over the same denominators and
  numerators.
\end{proof}

\begin{theorem}[The output under NO\_CHECK]\label{thm:noCheck}
  \uses{lem:findPoints_searchList, lem:noCheckList, thm:findPoints, def:IsPoint}
  Proposed names: \lname{Ratpoints.exists\_eq\_of\_mem\_of\_findPoints\_eq\_ok\_of\_noCheck},
  \lname{Ratpoints.mem\_of\_isPoint\_of\_findPoints\_eq\_ok\_of\_noCheck},
  \lname{Ratpoints.map\_sublist\_of\_findPoints\_eq\_ok\_of\_noCheck}. For well-formed parameters
  and options with $\mathrm{noCheck}$, if $\mathrm{findPoints}(c, H, P, o)$ returns the list $L$:
  every member of $L$ is $(x, 0, z)$ with $\gcd(x, z) = 1$, $z > 0$ or ($z = 0$ and $x = 1$), $|x|
  \le H$ and $z \le H$; for every normalized point $(x, y, z)$ of height at most $H$, $(x, 0, z)$
  lies in $L$; and if the search with $\mathrm{noCheck}$ cleared and $\mathrm{noY}$ set returns
  $L'$, then $L'$ with $y$ replaced by $0$ is a sublist of $L$, so that the points come in the order
  of the search among the survivors.
\end{theorem}

\begin{proof}
  Lemma~\ref{lem:findPoints_searchList} gives $L$ and $L'$ as the lists that the search delivers for
  one sound test, and Lemma~\ref{lem:noCheckList} the first and the third statement. For the second,
  the point $(x, |y|, z)$ lies in $L'$ by Theorem~\ref{thm:findPoints}, and the sublist puts $(x, 0,
  z)$ into $L$.
\end{proof}

No duplicates is the general statement of Theorem~\ref{thm:findPoints}, which holds for every
option: without $\mathrm{noCheck}$ the list is a permutation of the reference, filtered under
$\mathrm{noY}$, and under $\mathrm{noCheck}$ Lemma~\ref{lem:noCheckList} applies. These statements
say nothing of the survivors of a given set of moduli beyond their being normalized pairs: which
pairs pass the sieve is what the moduli decide, and the differential tests compare them with the C
program where both sieve with the same primes over the same range (Section~\ref{sec:outputTests}).

\section{The order}\label{sec:order}

The specification gains a file with the order of the points, \verb|Ratpoints/Spec/Order.lean|, which
uses Mathlib and nothing of the model and which the challenge restates verbatim. It imports
\verb|Mathlib.Data.Nat.Factors| for the prime factors of a number, in increasing order with
multiplicity.

\begin{definition}[The polynomial searched]\label{def:searchedPoly}
  \uses{def:binaryForm}
  \lean{Ratpoints.searchedPoly}
  \leanok
  For the switch $r$ of the reversal and the coefficient list $c$ of degree $d$, the polynomial that
  ratpoints searches and whether it is the reversed one:
  \[
    \mathrm{searchedPoly}(r, c) =
    \begin{cases}
      ((c_d, c_{d-1}, \dots, c_0), \mathrm{true}) & \text{if } r,\ d \text{ even and } (c_0 = 0
        \text{ or } (c_d \text{ a square},\ c_0 \text{ not})), \\
      ((0, c_d, c_{d-1}, \dots, c_1), \mathrm{true}) & \text{if } r,\ d \text{ odd},\
        |c_d| \ne 1,\ c_0 = 0,\ |c_1| = 1, \\
      (c, \mathrm{false}) & \text{otherwise},
    \end{cases}
  \]
  the three cases of the reversal check (Definition~\ref{def:reversal}) on the coefficient list. The
  binary form of a reversed list is that of $c$ with $x$ and $z$ exchanged, except for even $d$ with
  $c_0 = c_1 = 0$, which ratpoints refuses.
\end{definition}

\begin{definition}[The divisors and the denominators of the search]\label{def:binaryDivisors}
  \uses{def:binaryForm}
  \lean{Ratpoints.FactorsByTrialDivision, Ratpoints.binaryDivisors, Ratpoints.searchDenominators}
  \leanok
  A natural number $n$ is \emph{factored by trial division} when the product of its prime factors
  from $1021$ on, with multiplicity, is below $1021^2$, or none of its prime factors exceeds $1021$:
  ratpoints divides off $2$ and the odd primes up to $1021$ in turn and takes a cofactor below
  $p^2$, for the next prime $p$, as a prime, which it is, so that it fails exactly when the part of
  $n$ at the primes from $1021$ on is at least $1021^2$ and not a power of $1021$ (decidable, by an
  instance). The squarefree divisors of $n$ at most $H$ in the order of binary counting: for the
  distinct prime factors $p_1 < p_2 < \cdots$ of $n$, the products of the sublists of
  $(p_1, p_2, \dots)$ in the order of \verb|List.sublists|,
  $1, p_1, p_2, p_1 p_2, p_3, p_1 p_3, \dots$, with $p_1$ the lowest bit, those at most $H$. The
  denominators of the search of $c$, of degree $d$, for the height $H$ and the bound $m$ of the
  number of divisors: when $d$ is odd, $n = |c_d| \ne 1$, $n$ is factored by trial division and it
  has fewer than $m$ squarefree divisors at most $H$, the numbers $e k^2 \le H$, divisor $e$ by
  divisor in that order and $k = 1, 2, \dots$ for each; otherwise $1, \dots, H$. This is the order
  of the loop of the search, which visits those of them that its tests let through; the others carry
  no point.
\end{definition}

\begin{definition}[The points in the order of ratpoints]\label{def:orderedPoints}
  \uses{def:searchedPoly, def:binaryDivisors, def:IsPoint}
  \lean{Ratpoints.pointsWithDen, Ratpoints.orderedPoints}
  \leanok
  The points of height at most $H$ with the denominator $z$: for $x = -H, \dots, H$ in turn, with
  $y = \lfloor \sqrt{F_c(x, z)} \rfloor$, those among $(x, y, z)$ and $(x, -y, z)$ that are
  normalized points of height at most $H$, in this order, one triple when $y = 0$; the denominator
  $0$ gives the points at infinity, whose numerator is $1$. The points in the order of ratpoints,
  for the bound $m$ of the number of divisors, the switch $r$ of the reversal, $c$ and $H$: for
  $(c', \mathrm{rev}) = \mathrm{searchedPoly}(r, c)$, the concatenation of the points of $c'$ with
  the denominators $0$ and then those of the search of $c'$ in their order, each $(a, y, b)$ mapped
  to $(b, y, a)$ for $a \ge 0$ and to $(-b, y, -a)$ for $a < 0$ when $\mathrm{rev}$ is set, a point
  of the caller's curve.
\end{definition}

The file checks its definitions on outputs of the C program by \verb|decide +kernel|, which the
kernel decides and which, unlike \verb|native_decide|, adds no axiom:
\verb|ratpoints '-1 0 9' 10 -q|, which reverses the polynomial to \verb|9 0 -1|,
\verb|ratpoints '0 1 0 2' 10 -q|, which reverses it to \verb|0 2 0 1|,
\verb|ratpoints '-6 5 -3 30' 12 -q|, which visits the denominator $6$ before $5$, in the order $1,
2, 3, 6, 5, 10, 15, 30$ of the divisors of $30$, and \verb|ratpoints '0 1031 -11341 30930' 6 -q|,
whose leading coefficient $2 \cdot 3 \cdot 5 \cdot 1031$ is factored completely, so that $6$ comes
before $5$, while $30 \cdot 1042451$ of \verb|0 1042451 -11466961 31273530| is not and its
denominators ascend; $3 \cdot 1031$ is factored by trial division, its cofactor $1031 < 37^2$ taken
as a prime at $37$, and $3 \cdot 1042451$ and $1021 \cdot 1031$ are not. The bound of the number of
divisors is an argument of the specification, which the order theorem instantiates with the constant
$\mathrm{maxDivisors} = 512$ of the model (\verb|MAX_DIVISORS|), so that the build with other values
of the constants varies it ($2$ and $100000$).

\begin{lemma}[The rough part]\label{lem:roughPart}
  \uses{def:binaryDivisors}
  \lean{Nat.IsRoughPart, Nat.IsRoughPart.ne_zero, Nat.IsRoughPart.eq_prod_filter, Nat.IsRoughPart.succ, Nat.IsRoughPart.of_le_of_forall, Ratpoints.factorsByTrialDivision_of_isRoughPart_one, Ratpoints.factorsByTrialDivision_of_isRoughPart_of_lt, Ratpoints.eq_one_iff_factorsByTrialDivision_of_isRoughPart, Array.pairwise_map_toList_push}
  \leanok
  The part $w$ of $n$ at the primes from $q$ on, the \emph{rough part}: $n = k w$ with every prime
  factor of $k$ below $q$ (Mathlib's \verb|Nat.smoothNumbers|) and every prime factor of $w$ at
  least $q$, a notion for Mathlib. For $n \ne 0$ it is not $0$ and it is the product of the prime
  factors of $n$ from $q$ on, with multiplicity; dividing the rough part $p^v m$ at a prime $p$,
  with $p \nmid m$, by $p^v$ leaves the rough part at $p + 1$; and the rough part at $q$ is that at
  $q' \ge q$ when no prime lies between them. A number whose rough part at some $q \le 1022$ is $1$,
  or a number $n \ne 0$ whose rough part at some $q \le 1021$ is below $q^2$, is factored by trial
  division; and when the rough part of $n \ne 0$ at $1021$ is at least $1021^2$, $n$ is factored by
  trial division exactly when its rough part at $1022$ is $1$. The images of the members of an array
  after a push are in order exactly when those before it are and each of them comes before the image
  of the new member.
\end{lemma}

\begin{proof}
  \leanok
  From the definitions, with Mathlib's smooth numbers and the prime factors with multiplicity.
\end{proof}

\begin{lemma}[The divisors of the leading coefficient in binary counting]\label{lem:binaryDivisors}
  \uses{def:setupUs1, def:binaryDivisors}
  \lean{Ratpoints.Setup.foldl_extendDivisors_eq_filter, Ratpoints.Setup.IsFactorization.map_fst_eq_dedup, Ratpoints.Setup.foldl_extendDivisors_eq_binaryDivisors, Ratpoints.Setup.map_snd_run_setupUs1}
  \leanok
  For a list of positive numbers $p$, the fold of the step that appends $x p$ for the divisors
  $x \le \lfloor b_{\mathrm{high}}/p \rfloor$, from $(1)$, is the list of the products of the
  sublists of the list that are at most $b_{\mathrm{high}} \ge 1$, in the order of
  \verb|List.sublists|; the primes of a factorization of $n \ne 0$, in increasing order, are the
  distinct prime factors of $n$; so for the primes of a factorization of $|c_d|$ with the cofactor
  $1$ the fold is the list of the squarefree divisors of $|c_d|$ up to $b_{\mathrm{high}}$ in binary
  counting. So, for $c_d \ne 0$ and $1 \le b_{\mathrm{high}}$, the divisor set-up succeeds exactly
  when $|c_d|$ is factored by trial division and has fewer than $\mathrm{maxDivisors}$ squarefree
  divisors at most $b_{\mathrm{high}}$, and its divisors are then those of binary counting.
\end{lemma}

\begin{proof}
  \leanok
  By induction on the primes, the doubling step adding the products with the new prime after the old
  ones, as \verb|List.sublists| does; the rest from the contracts of the factorization and of the
  divisor set-up (Definition~\ref{def:setupUs1}).
\end{proof}

With the rough part, the contracts of the factorization and of the divisor set-up
(Definition~\ref{def:setupUs1}), \lname{Ratpoints.Setup.factorLcf\_spec} and
\lname{Ratpoints.Setup.setupUs1\_spec}, become one each, and that of the doubling loop,
\lname{Ratpoints.Setup.squarefreeDivisors\_spec}, covers the loop that stops: the factorization ends
with the cofactor $1$ exactly when the number is factored by trial division, and then with a
factorization in increasing order; the doubling loop gives the fold of its step or, when it stops at
$\mathrm{maxDivisors}$ entries, a prefix of it with at least $\mathrm{maxDivisors}$ entries; the
set-up returns its records and its divisors exactly when the factorization ends with the cofactor
$1$ and there are fewer than $\mathrm{maxDivisors}$ divisors, and nothing otherwise. They are proved
with one invariant of the loop of the factorization, which the node of Definition~\ref{def:setupUs1}
names with its lemmas: the product of the prime powers found and the cofactor is $|c_d|$, the primes
found ascend, and while primes of the table remain, the cofactor is the rough part of $|c_d|$ at the
next prime of the table, which needs that the table starts at $3$, ends at $1021$ and has no prime
between two neighbours; after the last prime, or when the cofactor is below the square of the next
prime, it is $1$ exactly when the factorization is complete. The invariants of the divisors keep,
once the loop has stopped at $\mathrm{maxDivisors}$ entries, that the list is a prefix of the full
fold.

\begin{lemma}[The denominators and the points of the search]\label{lem:searchDenominators}
  \uses{lem:binaryDivisors, def:searchDens, def:binaryDivisors, def:searchedPoly, def:orderedPoints,
    def:reversal, lem:shape_dens, def:bruteForce}
  Proposed names: \lname{Ratpoints.searchDens\_eq\_searchDenominators},
  \lname{Ratpoints.searchedPoly\_eq\_of\_reversal\_eq\_ok},
  \lname{Ratpoints.pointsOfDen\_eq\_pointsWithDen},
  \lname{Ratpoints.pointsOfDen\_zero\_eq\_pointsWithDen},
  \lname{Ratpoints.pos\_of\_mem\_binaryDivisors},
  \lname{Ratpoints.le\_and\_le\_of\_mem\_searchDenominators},
  \lname{Ratpoints.Sieve.squares1Dens\_one\_eq\_filter},
  \lname{Ratpoints.Setup.isSquare\_eq\_true\_iff}, \lname{Ratpoints.Setup.isSquare\_eq\_false\_iff},
  \lname{Ratpoints.BruteForce.pointsAbove\_eq\_ite}, \lname{Ratpoints.isPoint\_iff\_of\_le},
  \lname{Ratpoints.ite\_pointsAbove\_eq\_filter}, \lname{List.filter\_eq\_ite\_of\_forall}. For $1
  \le H$, the denominators of the search of $c$ of its degree (Definition~\ref{def:searchDens}) are
  the denominators of the search of the specification with the bound $\mathrm{maxDivisors}$; a
  reversal check that does not refuse $c$ hands on the polynomial searched of the specification with
  its flag; and what the brute-force reference finds over a denominator $0 < z \le H$ and the
  numerators $-H, \dots, H$ are the points with the denominator $z$, and over the candidate $(1, 0)$
  those with the denominator $0$.
\end{lemma}

\begin{proof}
  \uses{lem:flatMap_sound}
  Lemma~\ref{lem:binaryDivisors} for the divisor set-up, and the divisors times squares of the shape
  as a filter of the products $e k^2$; the reference checks a candidate by its square root and its
  test for common factors, which on the box are the conditions of a point.
\end{proof}

\begin{theorem}[The order of the points]\label{thm:order}
  \uses{thm:findPoints_order, lem:searchDenominators, def:orderedPoints, def:findPoints}
  Proposed name: \lname{Ratpoints.eq\_orderedPoints\_of\_findPoints\_eq\_ok}. For well-formed
  parameters and options with $\mathrm{noCheck}$ false, if $\mathrm{findPoints}(c, H, P, o)$ returns
  the list $L$, then $L$ is the list of the points in the order of ratpoints for the bound
  $\mathrm{maxDivisors}$, the switch of the reversal of $o$, $c$ and $H$, without the points with $y
  < 0$ when $\mathrm{noY}$ is set.
\end{theorem}

\begin{proof}
  Theorem~\ref{thm:findPoints_order} gives $L$ as what the reference of the polynomial searched
  finds over the candidate $(1, 0)$ and the denominators of its search, mapped back after a reversal
  and filtered by $\mathrm{noY}$, and Lemma~\ref{lem:searchDenominators} turns each part into its
  form in the specification.
\end{proof}

\section{What comparator certifies}\label{sec:certifiedG}

\verb|comparator/Challenge.lean| restates, besides $\mathrm{degree}$, $\mathrm{binaryForm}$,
$\mathrm{binaryFormPoly}$ and $\mathrm{IsPoint}$, the decidability instance of $\mathrm{IsPoint}$,
which the points with a denominator filter by, and the definitions of
\verb|Ratpoints/Spec/Order.lean| with their instance, verbatim, and imports exactly the public
imports of the specification files it restates, \verb|Mathlib.Data.Nat.Factors| joining those of
\verb|Ratpoints/Spec/Basic.lean| with the file of the order; from the model it takes
\verb|bruteForce|, \verb|findPoints|, \verb|findPointsWith|, \verb|Answer|, \verb|Delivery|,
\verb|Delivery.deliver|, \verb|Options|, \verb|Params|, \verb|Params.WF|, \verb|SieveOptions|,
\verb|sieveParams|, \verb|maxDegree| and \verb|maxDivisors|. The statements come in with their
theorems: the protocol with G.7, the membership statement with the options of the output with G.8,
the three statements under \verb|RATPOINTS_NO_CHECK| with G.9, and the order theorem with the
restated file of the order with G.10, each item running comparator; G.14 runs it with
\verb|--paranoid| (Section~\ref{sec:outputFiles}). The twelve statements, for all $c$, $H$,
parameters $P$ and options $o$:
\begin{enumerate}
  \item a triple lies in $\mathrm{bruteForce}(c, H)$ exactly when it is a normalized point of height
    at most $H$ (\verb|challenge_mem_bruteForce_iff|, unchanged);
  \item $\mathrm{bruteForce}(c, H)$ has no duplicates (\verb|challenge_nodup_bruteForce|,
    unchanged);
  \item for $P$ well formed, $\mathrm{noCheck}$ false and $\mathrm{findPoints}(c, H, P, o) = L$: $p
    \in L$ exactly when $p$ is a normalized point of height at most $H$ with $y \ge 0$ if
    $\mathrm{noY}$ is set (\verb|challenge_mem_iff_isPoint_of_findPoints_eq_ok|, with the options of
    the output, G.8);
  \item for $P$ well formed and $\mathrm{findPoints}(c, H, P, o) = L$, $L$ has no duplicates
    (\verb|challenge_nodup_of_findPoints_eq_ok|, its text unchanged, now over the options of the
    output too);
  \item with the Sturm step on, $\mathrm{findPoints}$ refuses exactly when $c$ has fewer than two or
    more than $\mathrm{maxDegree} + 1$ entries, $H = 0$, $H \ge 2^{63}$ or $\Phi_c$ is not
    squarefree (unchanged);
  \item with the Sturm step off, exactly when $c$ has fewer than two or more than
    $\mathrm{maxDegree} + 1$ entries, $H = 0$, $H \ge 2^{63}$, $d + 2 \le N$, or the reversal is
    allowed, $d$ is even and $c_0 = c_1 = 0$ (unchanged);
  \item the parameters of every set of options of the program are well formed
    (\verb|challenge_wf_sieveParams|, unchanged);
  \item for $P$ well formed, $\mathrm{noCheck}$ false and $\mathrm{findPoints}(c, H, P, o) = L$: $L$
    is the list of the points in the order of ratpoints for $\mathrm{maxDivisors}$, the switch of
    the reversal of $o$, $c$ and $H$, filtered by $\neg\mathrm{noY} \lor y \ge 0$
    (\verb|challenge_eq_orderedPoints_of_findPoints_eq_ok|, new, G.10);
  \item[9--11.] for $P$ well formed, $\mathrm{noCheck}$ set and $\mathrm{findPoints}(c, H, P, o) =
    L$: every member of $L$ is $(x, 0, z)$ with $\gcd(x, z) = 1$, $z > 0$ or ($z = 0$ and $x = 1$),
    $|x| \le H$ and $z \le H$; for every normalized point $(x, y, z)$ of height at most $H$, $(x, 0,
    z) \in L$; and if the search with $\mathrm{noCheck}$ cleared and $\mathrm{noY}$ set returns
    $L'$, the list $L'$ with $y$ replaced by $0$ is a sublist of $L$
    (\verb|challenge_exists_eq_of_mem_of_findPoints_eq_ok_of_noCheck|,
    \verb|challenge_mem_of_isPoint_of_findPoints_eq_ok_of_noCheck|,
    \verb|challenge_map_sublist_of_findPoints_eq_ok_of_noCheck|, new, G.9);
  \item[12.] for $P$ well formed and every callback with its data, $\mathrm{findPointsWith}$ with
    them gives, mapped to its data and count, what the delivery of the list of
    $\mathrm{findPoints}(c, H, P, o)$ to the callback from its data, the count $0$ and the flag
    clear gives, and the same refusal, with the same error, exactly where $\mathrm{findPoints}$
    refuses (\verb|challenge_map_findPointsWith_eq_map_findPoints|, new, G.7).
\end{enumerate}
The protocol speaks about the reference of the delivery of the model, a fold of a few lines, which
one reads with the challenge as one reads the specification; \verb|findPointsWith| is the code that
the executable runs, with the callback of \verb|main.c|; \verb|findPoints|, of which statements 3 to
6 and 8 to 11 speak, is its instance with the appending callback (Definition~\ref{def:findPoints}),
and statement 12 links the two. The theorems of the library behind the twelve statements have guards
on their axioms, which are the standard ones.

\section{The program}\label{sec:outputCli}

The program reads the options of the output and runs the search with the callback of \verb|main.c|,
which prints the points as the search delivers them.

\begin{definition}[The settings of the output]\label{def:cliOutput}
  \uses{def:Options}
  Proposed names: \lname{Ratpoints.Cli.Output}, \lname{Ratpoints.Cli.Output.store},
  \lname{Ratpoints.Cli.OptionArg.word}, \lname{Ratpoints.Cli.defaultFormat}. The global variables of
  \verb|main.c| with their defaults, those of a command line without options: whether the report is
  left out (\verb|-q|), whether it is that of \verb|-v| (the bit \verb|RATPOINTS_VERBOSE| of the
  flags, \csource{ratpoints.h}{407}{407}), whether one point is enough (\verb|-1|), whether the
  points at infinity are kept (\verb|-I|, and not \verb|-i|), whether no point is printed
  (\verb|-z|, and not \verb|-Z|), the format of a point (\verb|-f|), none for the default, and the
  strings before, between and after the points (\verb|-fs|, \verb|-fm|, \verb|-fe|), none for none
  (\csource{main.c}{57}{64}, \csource{main.c}{383}{390}); the setting that an option changes; the
  kind of argument of \verb|-f|, \verb|-fs|, \verb|-fm| and \verb|-fe|, the next word, whatever it
  is, also one that starts with \verb|-|; and the format that \verb|main.c| takes when the command
  line gives none, \verb|(%x : %z)\n| when the points are not checked exactly or only one point of
  each pair is delivered, and \verb|(%x : %y : %z)\n| otherwise (\csource{main.c}{681}{685}, which
  leaves the line of a point of stage~B).
\end{definition}

The reading of the options (Definition~\ref{def:cliInput}) reads \verb|-x -X -y -Y| into
$\mathrm{noCheck}$ and $\mathrm{noY}$ of the options of the search and
\verb|-q -v -z -Z -1 -i -I -f -fs -fm -fe| into the settings of the output, of each the last one
winning; \verb|-f| followed by a character other than \verb|s|, \verb|m| and \verb|e|, and a missing
word, are error~6; after the options the format is the default unless \verb|-f| gave one, and
\verb|-q| cancels \verb|-v|. Its C ranges gain \csource{main.c}{533}{577},
\csource{main.c}{602}{637} and \csource{main.c}{688}{688}, the head and the end of \verb|read_input|
(\csource{main.c}{315}{316}, \csource{main.c}{691}{692}), which go to \verb|Ratpoints.Cli.input|,
and the array of the coefficients (\csource{main.c}{54}{54}), which goes to
\verb|Ratpoints.Cli.coefficients|; and only the options of the search region and the range of the
denominators, \verb|-l -u -dl -du|, keep the exit status 12 (stage~H).

\begin{definition}[The formats and the strings]\label{def:formatPoint}
  Proposed names: \lname{Ratpoints.Cli.escape}, \lname{Ratpoints.Cli.printString},
  \lname{Ratpoints.Cli.formatPoint}. What \verb|main.c| prints for a backslash followed by a
  character in a format or a string: a tab for \verb|t|, a line feed for \verb|n|, a backslash for a
  backslash, a percent sign for \verb|%|, and the backslash and the character as they are otherwise
  (\csource{main.c}{178}{182}, \csource{main.c}{961}{965}, the two switches of \verb|process| and
  \verb|print_string| as one function). What \verb|print_string| prints for a string: nothing for
  none, otherwise its characters, a backslash with the character after it as above and a backslash
  at the end as itself, the percent sign not special (\csource{main.c}{952}{960},
  \csource{main.c}{966}{968}). What \verb|process| prints for the point $(x : y : z)$ in a format,
  character by character: \verb|%x|, \verb|%y| and \verb|%z| the coordinates in decimal, but
  \verb|%y| itself when the points are not checked exactly or only one point of each pair is
  delivered, a percent sign before another character both, a percent sign or a backslash at the end
  itself, a backslash and the character after it as above, and every other character itself
  (\csource{main.c}{157}{177}, \csource{main.c}{183}{186}).
\end{definition}

\begin{definition}[The callback of the program]\label{def:cliProcess}
  \uses{def:delivery, def:formatPoint, def:cliOutput}
  Proposed names: \lname{Ratpoints.Cli.Data}, \lname{Ratpoints.Cli.Data.ofOutput},
  \lname{Ratpoints.Cli.process}. The data that \verb|main.c| passes to \verb|process| through
  \verb|info|: whether a point was printed, whether no point is printed, whether one point is
  enough, whether the points at infinity are kept, whether \verb|%y| prints itself, the format, the
  strings before and between the points, and the text printed so far (\csource{main.c}{110}{113});
  its value for the settings of the output and the options of the search, no point printed yet,
  \verb|%y| printing itself under either option of the output (\csource{main.c}{235}{242}). The
  callback on $(x, z, y)$: a point at infinity is passed over when the points at infinity are not
  kept, with the weight $0$ and the flag clear; every other point has the weight $1$ and stops the
  search when one point is enough; it is printed unless no point is printed, after the string before
  the points if it is the first point printed and after the string between them otherwise, in the
  format (\csource{main.c}{139}{142}, \csource{main.c}{153}{156}, \csource{main.c}{187}{189}; the
  test of the search region, lines 150-152, is stage~H's).
\end{definition}

\begin{definition}[The report]\label{def:cliReport}
  \uses{def:cliOutput, def:Outcome, def:cliFailure, def:sieveParams}
  Proposed names: \lname{Ratpoints.Cli.printPoly}, \lname{Ratpoints.Cli.Message},
  \lname{Ratpoints.Cli.Message.text}, \lname{Ratpoints.Cli.preamble}, \lname{Ratpoints.Cli.report}.
  The polynomial as \verb|print_poly| prints it: its terms from the highest degree down, the zero
  ones left out, the first after \verb|- | if it is negative and the others after \verb| + | or
  \verb| - |, each the absolute value of its coefficient, followed for the degree $1$ by a space and
  \verb|x| and for a degree $i \ge 2$ by a space and \verb|x^i| (the constant term is the number
  alone), the coefficient and the space left out when it is $1$ there, and a line feed
  (\csource{main.c}{913}{949}). The messages of the report that the program prints, with their data
  and their texts: message 0, the banner of \verb|ratpoints-lean|; message 2, the count of the
  points found, of the pairs under either option of the output; message 5, the curve; message 6, the
  height bound; and message 7, what the search did, from the flags written back, the Jacobi switch
  and $\mathrm{noCheck}$ (\csource{main.c}{972}{976}, \csource{main.c}{979}{983},
  \csource{main.c}{999}{1021}, \csource{main.c}{1031}{1033}). What \verb|main.c| prints before the
  search: nothing under \verb|-q|, otherwise the banner, the curve and, unless \verb|-v|, the height
  bound and a blank line (\csource{main.c}{215}{220}); and after it: nothing under \verb|-q|,
  otherwise, unless \verb|-v|, two line feeds and message 7, then message 2 with the count of the
  search (\csource{main.c}{268}{276}).
\end{definition}

The program (Definition~\ref{def:cliRun}, \lname{Ratpoints.Cli.run}) reads its command line, prints
the head of the report, and runs $\mathrm{findPointsWith}$ with the parameters
$\mathrm{sieveParams}(s)$, the options of the search and the callback of the program on its data; on
a refusal it prints the message of the failure and exits with its status, and otherwise it prints
what the callback printed, the string before the points if it printed no point, the string after
them and the end of the report, and exits with status~$0$ (\csource{main.c}{195}{198},
\csource{main.c}{251}{252}, \csource{main.c}{265}{265}, \csource{main.c}{267}{267},
\csource{main.c}{284}{285}, the head and the end of \verb|main| joining the links it has). The text
is written once, at the end.

\begin{theorem}[What the program prints]\label{thm:cliRun}
  \uses{def:cliProcess, def:cliReport, thm:protocol, def:cliRun, def:cliInput, def:sieveParams}
  Proposed names: \lname{Ratpoints.Cli.countedPoints}, \lname{Ratpoints.Cli.Data.printPoint},
  \lname{Ratpoints.Cli.call\_process\_of\_keep}, \lname{Ratpoints.Cli.call\_process\_of\_not\_keep},
  \lname{Ratpoints.Cli.deliver\_process}, \lname{Ratpoints.Cli.total\_deliver\_process},
  \lname{Ratpoints.Cli.out\_append\_deliver\_process\_ofOutput},
  \lname{Ratpoints.Cli.info\_total\_eq\_of\_findPointsWith\_eq\_ok},
  \lname{Ratpoints.Cli.exists\_findPointsWith\_eq\_ok\_of\_findPoints\_eq\_ok},
  \lname{Ratpoints.Cli.run\_eq\_of\_findPointsWith\_eq\_ok},
  \lname{Ratpoints.Cli.run\_eq\_of\_findPoints\_eq\_ok},
  \lname{Ratpoints.Cli.run\_eq\_of\_onePoint},
  \lname{Ratpoints.Cli.total\_eq\_of\_findPointsWith\_eq\_ok},
  \lname{Ratpoints.Cli.preamble\_eq\_empty\_of\_quiet},
  \lname{Ratpoints.Cli.report\_eq\_empty\_of\_quiet},
  \lname{String.intercalate\_cons\_eq\_append\_join}. The points of a list that the callback counts:
  those that \verb|-i| keeps, every point with $z \ne 0$ and the points at infinity unless
  \verb|-i|, and of them only the first under \verb|-1|. If the command line reads as the
  coefficients $c$, the height $H$, the values $s$ of the options of the parameters, the options $o$
  of the search and the settings $u$ of the output, and $\mathrm{findPoints}(c, H,
  \mathrm{sieveParams}(s), o)$ returns $L$, the program exits with status~$0$ and prints the head of
  the report, the string before the points, the formats of the counted points of $L$ joined by the
  string between them (none under \verb|-z|), the string after them, and the end of the report,
  whose count is the number of the counted points; under \verb|-q| only the strings and the points;
  and under \verb|-q| and \verb|-1| without \verb|-z|, when $p$ is the first point of $L$ that
  \verb|-i| keeps, the string before the points, the format of $p$ and the string after them:
  \verb|-1| prints the first point of the list that \verb|-i| keeps.
\end{theorem}

\begin{proof}
  The delivery of any list to the callback of the program prints the counted points of the list in
  turn, counts them, and sets the flag exactly under \verb|-1| when a point is kept; by the protocol
  of the delivery (Theorem~\ref{thm:protocol}), with the well-formed parameters of the program, the
  search with the callback gives the data and the count of the delivery of the list of
  $\mathrm{findPoints}$.
\end{proof}

\section{The differential tests}\label{sec:outputTests}

The harness compares what \verb|ratpoints-lean| prints with the C program at the pin and with its
recorded outputs, now with the options of the output and the report:
\begin{itemize}
  \item an output of points in the default format with \verb|-q| is compared as a list of points
    with a recorded output (\verb|test4|, \verb|rptest|), and exactly, byte for byte, with an output
    of the C program, as is an output with the options of the output (\verb|-x -y -f -fs -fm -fe|),
    also one that does not end with a line feed;
  \item a report, a command line without \verb|-q|, is compared exactly, with the banner exchanged
    for that of the C program and the messages that \verb|ratpoints-lean| does not print dropped
    from the report of the C program only: those on the search region and the search intervals
    (messages 3 and 8, stage~H) and on the moduli used (message 4, stage~I), as the filter
    \verb|m()| of \verb|test4.sh| drops the last;
  \item under \verb|-v|, where the C program prints the report of its library too,
    \verb|ratpoints-lean| is compared with the report of the C program on the same command line
    without \verb|-v|, its height and its messages 4, 7 and 8 dropped, as \verb|main.c| does under
    \verb|-v|, and the lines of that text must be among those of the output of the C program with
    \verb|-v|, in their order, with the same lines of points and pairs (with a format or a string,
    whose points need not end their lines, by characters);
  \item the outputs that \verb|test4.sh| filters through \verb|grep found|,
    \verb|grep 'Curve equation'| and \verb|grep -c .| are compared after the same filter, the banner
    exchanged first, with the exit status~$0$;
  \item under \verb|-x|, whose lines are the survivors of the sieve, which depend on the moduli and
    on the search region: when pinning the moduli (\verb|-p 30 -N 30| appended, which keeps the
    moduli of the model without options, the first $30$ odd primes, but puts all of them into the
    first two stages, with no third stage, where the model without options puts the last $8$ into
    its third stage) does not change the output of \verb|ratpoints-lean|, exactly against the C
    program on the command line with \verb|-p 30 -N 30 -s| appended, which then sieves with the same
    primes over the whole range of the height, without the region of its Sturm step, as the model
    does until stages~H and~I; and, without a format or a string, also by a weaker check against the
    C program on the command line itself, the only check when the pin changes the output of
    \verb|ratpoints-lean|: every line $(x : z)$ a normalized coprime pair in the range of the
    height, none repeated, the pairs that are points the same and in the same order in both outputs,
    the pairs common to both in the same order, the pairs of each output in one block at the same
    place among the other lines, under \verb|-1| one line in each output when one of them is a point
    (the same if both are) and else at most one, the other lines equal, and a count line equal but
    for its number, which must be that of the pairs printed (unless \verb|-z|);
  \item exit codes are compared exactly, and \verb|ratpoints-lean| must write nothing to its
    standard error: it runs with \verb|LEAN_ABORT_ON_PANIC=1|, so that a panic of the Lean runtime
    ends it;
  \item two searches at the height $10^{12}$ that both programs stop at a point at infinity before
    the sieve, \verb|'1 0 1'| and \verb|'3 2 -1 12'| with \verb|-q -1|, are run with a timeout of
    $30$ seconds and a cap of $4$~GB on the memory;
  \item a new default suite \verb|outputs| runs $300$ random polynomials at small heights with
    random options of the output, \verb|-q| or \verb|-v|, formats and strings, pinned at times,
    against the C program; the errors suite has the cases of the options of the output; the four
    runs of \verb|test4.sh| whose \verb|-f|, \verb|-fs|, \verb|-fm| or \verb|-fe| lacks its word
    fail with error~6, as in the C program; and the fuzz draws the options of the output and format
    words, kept in UTF-8, on an eighth line of each round;
  \item a command line with \verb|-l|, \verb|-u|, \verb|-dl| or \verb|-du| must fail as the C
    program does if that fails before the option, and else exit with status 12, also under the
    filters of \verb|test4.sh|.
\end{itemize}
Skipped are the outputs that \verb|test4.sh| (four runs) and \verb|test3.sh| (three runs) filter to
look at the report of the library under \verb|-v|; the two runs of \verb|test3.sh| whose filter
\verb|sed -n '/moduli used/,$p'| keeps the report from message 4 on, which needs messages 4
(stage~I) and 8 (stage~H); searches under \verb|-x| with a format or a string whose output the pin
changes; and searches above the largest height of the tests, $30000$, which \verb|test4high|
compares. In the prototype the default suites compare $3784$ times with none failed: \verb|test4|
$565$ (four outputs filtered, $13$ runs left to \verb|test4high| for their height), \verb|test3|
$16$ (five outputs filtered), \verb|rptest| $1206$, \verb|random| $300$, \verb|outputs| $300$, the
errors suite $313$ (the searches stopped at $10^{12}$ among them), the fuzz $800$, the parameters
$272$ and the high suite $12$, with $67$ runs under \verb|-x| compared exactly, none by the weaker
check alone; and \verb|test4high| compares $13$, none failed. Every point list is compared in its
order.

The \verb|#guard| examples of the model check the delivery against outputs of the C program and
against the reference: a callback that stops at its $k$-th call, with the weight $2$, for every $k$
up to beyond the length of a list; the stop and the flags in the exact check and at infinity; a stop
at the $k$-th call through one denominator, and $\mathrm{noY}$ and $\mathrm{noCheck}$ through it;
the prefix and the weights of a stop through the search with a callback; a stop at infinity at the
heights $10^{12}$ and $2^{63} - 1$; the flags written back against message 7 of the C program for
six curves, reversed, of odd degree with $|c_d| = 1$ and with divisors among them; and
$\mathrm{noY}$ and $\mathrm{noCheck}$ against runs of \verb|test4.sh|. Those of the program check
\verb|-1|, \verb|-1 -i|, \verb|-y| and \verb|-z -1|, the formats and the strings with their escapes,
a trailing \verb|%| and \verb|%y| under \verb|-y|, \verb|-i| after a reversal, whole reports under
\verb|-z|, under \verb|-v -1 -y| and with error 8, the curves of \verb|print_poly|, the lines of
message 7, and the reading of the options, the last setting winning, \verb|-f -x| taking \verb|-x|
as the format and a missing word as error 6, each against the C program at the pin up to the banner
and the messages that the program leaves out.

\section{Deviations}\label{sec:outputDeviations}

Each component differs from the C code in a stated way, recorded in its \verb|Deviation:| paragraph:
\begin{itemize}
  \item the callback a function that returns what the call gives, the weight, the flag and the data,
    for the \verb|int| that \verb|process| returns, the \verb|int| that it writes through
    \verb|quit| and the data behind \verb|info|, with no argument for \verb|quit|, which is clear at
    every call; one value for the delivery, the data, one running count and one flag, passed and
    returned, for the pointers \verb|info| and \verb|quit| (the local flag of
    \verb|find_points_work_1| and the field \verb|quit| of the sieving thread) and the counts
    \verb|total| that each function of the search keeps from zero and returns, which its caller adds
    to its own; the three branches of the map of the reversal at each call of the exact check one
    definition;
  \item the count an integer with no range obligation, since the weights that a callback returns are
    arbitrary \verb|int|s (the library sums them in a \verb|long|, which a caller whose weights sum
    to a negative value cannot tell from the codes of an error); for the callback of the program, of
    the weights $0$ and $1$, the count is at most the number of calls;
  \item the walk returns after a call that stops the search and its callers, the arms of the second
    stage and the scan, test the flag after it, where the C code returns from
    \verb|_ratpoints_sift0| from within the walk; one test of the flag after each denominator in the
    run, for the returns of \verb|run_block| in each shape and the end of the loop over the blocks;
  \item the two arms of the second stage functions of their own, \verb|siftBitArray| choosing
    between them by its argument;
  \item the sieving a function of its own for the block ``now do the sieving'', which takes over the
    deviations of the search on its lines (the denominators $1, \dots, H$, the row offsets made with
    zeros);
  \item the denominators of a shape listed before the first is tested and sifted, all of them, also
    those after a stop, which the loops of \verb|run_block| never reach (the blocks, stage~I): a
    stop in the sieve does not shorten the listing, so that a stop at a large height costs time and
    memory linear in the height, while a stop at infinity returns before the listing;
  \item the bits of the flags that the exact check and the points at infinity read as a structure of
    three Booleans, copied from the options of the search and the reversal check, for the bits of
    \verb|args->flags|; the options of the output as fields of the options of the search, as the
    switch of the reversal is;
  \item the result with the data, the count and three bits of the flags, the refusals returned
    apart, for the \verb|long| that \verb|find_points| returns, which carries the codes of the
    errors too, and the bits that it writes into \verb|args->flags|; nothing else written back:
    neither the reversed coefficients and the degree nor the bits \verb|RATPOINTS_CHECK_DENOM|,
    \verb|RATPOINTS_USE_JACOBI| and \verb|RATPOINTS_NO_REVERSE_AUTO|, which the program does not
    read; the search intervals (\verb|num_inter|, \verb|domain|, read by message 8) left to stage~H,
    and the numbers of the moduli used to stage~I;
  \item no failure for lack of memory, which the model cannot meet: neither a failed allocation of
    the sieving and the search (\verb|RATPOINTS_NO_MEMORY|) nor error 10 of the copies of the
    strings of the output that \verb|main.c| allocates;
  \item under \verb|RATPOINTS_NO_CHECK| the model delivers the survivors of all numerators $-H,
    \dots, H$, where the C code delivers only those in the intervals that its Sturm step keeps, none
    when that region misses $[-H, H]$ (\verb|sturm_empty|, only the points at infinity), and nothing
    at all when the form is negative everywhere (line 4889, even degree, no point at infinity): the
    search region is stage~H's, so that only there the depth of the Sturm step changes what is
    delivered;
  \item in the program, the data of \verb|main.c| a value, the callback's state, with the text
    printed so far, for the pointer \verb|info| and the standard output, and the output written once
    at the end, where \verb|main.c| prints and flushes after each point; the messages 4 (the moduli
    used, stage~I) and 8 (the search intervals, stage~H) left out, and message 3 (the search
    region), which needs \verb|-l| or \verb|-u|; the report of the library under \verb|-v|, its
    instrumentation, not printed; the strings of \verb|-f|, \verb|-fs|, \verb|-fm| and \verb|-fe|
    kept as values, for the copies that \verb|main.c| allocates;
  \item in the output, one function for the two switches of the escapes of \verb|process| and
    \verb|print_string|; recursions over the characters, with an integer for the GMP integer $y$,
    and a fold over the coefficients, for the loops over C strings and over the coefficients of
    \verb|print_poly|; an optional string, none for the null pointer, for a string of the output;
    and the bit \verb|RATPOINTS_VERBOSE| a field of the settings of the output;
  \item the formats and the strings read as Lean strings, which the runtime decodes from the command
    line as UTF-8, a byte sequence that is not UTF-8 read as U+FFFD, where \verb|main.c| copies the
    bytes; every escape of \verb|main.c| is ASCII and no byte of a multibyte character of UTF-8 is,
    so that the bytes printed agree on every format word in UTF-8;
  \item \verb|-l|, \verb|-u|, \verb|-dl| and \verb|-du| answered with the exit status 12 until
    stage~H.
\end{itemize}
The points at infinity follow the C code without a test $s \ne 0$, and the search writes back the
three bits of the flags that the program reads, on every path, so neither is a deviation. The lines
of this part of the C code that are neither modelled nor left to a later stage are excluded in
\verb|correspondence/ratpoints.toml|, each by the item named in parentheses: the prototypes of
\verb|find_points_init| (\csource{ratpoints.h}{439}{439}), interface, and of
\verb|find_points_clear| (\csource{ratpoints.h}{444}{444}), memory, as their definitions are (G.7);
in \verb|main.c| (G.11) the includes (\csource{main.c}{29}{37}) and the prototypes of its functions
(\csource{main.c}{94}{103}), other, since the model has neither headers nor forward declarations;
the GMP integers \verb|tmp| and \verb|tmp2| (\csource{main.c}{55}{55}) and their initialization and
release in \verb|main| with those of the coefficients and of the ends of the intervals
(\csource{main.c}{201}{207}, \csource{main.c}{279}{283}), memory; the declaration of the arguments
of the library (\csource{main.c}{199}{199}) and the initialization of the arguments that the model
has no counterpart for (\csource{main.c}{330}{332}), interface; the defaults of the tuning options
that the model checks and ignores (\csource{main.c}{339}{343}, \csource{main.c}{345}{346}),
parameters; and (G.12) the text of message 1 (\csource{main.c}{977}{978}), other, which \verb|main|
prints only when \verb|read_input| returns a value other than $0$, which it never does: its only
return is at line 691, and every other failure exits through \verb|error|. Left free for stage~H are
the lines of the search region of \verb|main.c| (among them the test of \verb|process|, lines
150-152, the messages 3 and 8, and the fields of \verb|args| that \verb|read_input| sets for it) and
the flag \verb|RATPOINTS_NO_REVERSE_AUTO| (\verb|ratpoints.h:421|), which only line 4687 of the
search region sets; and for stage~I message 4 and the numbers of the moduli used, zeroed on entry
with \verb|sp3| (\csource{find_points.c}{4551}{4552}) and set after the sieve
(\csource{find_points.c}{5163}{5163}, \csource{find_points.c}{5212}{5213}).

\section{What changes in the earlier chapters}\label{sec:outputEarlier}

Each change is made by the item of Section~\ref{sec:outputFiles} named in parentheses, so that the
blueprint stays true between the commits; the new uses are those that the dependency audit asks for
once the names have moved as stated.
\begin{itemize}
  \item Chapter~\ref{chap:spec}: the paragraph after Definition~\ref{def:bruteForce} says that the
    order is spelled out in Theorem~\ref{thm:order} and the protocol in Theorem~\ref{thm:protocol}
    (G.10); Section~\ref{sec:comparator-ab} states membership since stage~G with $\mathrm{noCheck}$
    false (G.8) and the statements of stage~G and the restated file of the order, whose public
    import the challenge takes (each with its item; the count and the record with G.14); the list of
    Section~\ref{sec:cli}: the program reads the options of the output and answers only
    \verb|-l -u -dl -du| with the status 12 (G.11) and prints the report without \verb|-q|, but for
    messages 3, 4 and 8 (G.12); the command line (Definition~\ref{def:cliInput}) reads the options
    of the output into the options of the search and the settings of the output, its node gaining
    \csource{main.c}{533}{577}, \csource{main.c}{602}{637} and \csource{main.c}{688}{688}, the head
    and the end of \verb|read_input| (\csource{main.c}{315}{316}, \csource{main.c}{691}{692}) and
    the array of the coefficients (\csource{main.c}{54}{54}), and using
    Definition~\ref{def:cliOutput} (G.11); the program (Definition~\ref{def:cliRun}) loses
    \verb|Ratpoints.Cli.pointLine|, whose range goes to Definition~\ref{def:cliOutput}, runs
    $\mathrm{findPointsWith}$ with the callback of the program and prints the report, with the links
    \csource{main.c}{195}{198}, \csource{main.c}{251}{252}, \csource{main.c}{265}{265},
    \csource{main.c}{267}{267} and \csource{main.c}{284}{285}, and uses
    Definitions~\ref{def:findPointsWith}, \ref{def:cliProcess} and~\ref{def:cliReport} in place of
    Definition~\ref{def:findPoints} (G.11, the report with G.12); the section on the differential
    tests: its list of what the later stages added gains stage~G, the suite \verb|outputs|, the
    options of the output and the reports compared exactly (Section~\ref{sec:outputTests}; G.11, the
    reports with G.12); the run of \verb|test4.sh| compares $565$ of its $582$ runs, the outputs
    filtered through \verb|grep found|, \verb|grep 'Curve equation'| and \verb|grep -c .| after the
    same filter with the banner exchanged, only the four filtered runs that look at the report of
    the library skipped and $13$ left to \verb|test4high| (G.12); the hand-written command lines,
    their number as at the end of the stage, also cover the options of the output, the formats and
    the strings of \verb|-f -fs -fm -fe| and the report, and the two searches at the height
    $10^{12}$ that both programs stop at a point at infinity, run with a timeout and a cap on the
    memory (G.11, the report with G.12); the fuzz runs $800$ random command lines, eight in each
    round, also drawing the options of the output and format words kept in UTF-8 (G.11); and the
    closing paragraph: a report is compared exactly, the banner exchanged for that of the C program
    and the messages 3, 4 and 8 dropped from the report of the C program only, the refusal of a form
    that is not squarefree without \verb|-q| compared as a whole report, and the output under
    \verb|-v| by its lines (G.12), \verb|ratpoints-lean| writes nothing to its standard error
    (\verb|LEAN_ABORT_ON_PANIC=1|), and the options without support are \verb|-l -u -dl -du| (G.11),
    the counts taken from the run of the last item; the files section names
    \verb|Ratpoints/Spec/Order.lean| (G.2).
  \item Chapter~\ref{chap:sieve}: the parameters (Definition~\ref{def:Params}) gain
    \csource{ratpoints.h}{406}{406}, the bit \verb|RATPOINTS_NO_JACOBI| for which their Jacobi
    switch stands (G.8); the run (Definition~\ref{def:sieveRun}) delivers to the callback and
    returns after a denominator at which the callback stopped the search, its node naming
    \lname{Ratpoints.Sieve.SieveRunInvariant} as a structure with
    \lname{Ratpoints.Sieve.SieveRunInvariant.quit\_eq},
    \lname{Ratpoints.Sieve.SieveRunInvariant.eq\_deliver},
    \lname{Ratpoints.Sieve.SieveRunInvariant.eq\_deliver\_append},
    \lname{Ratpoints.Sieve.SieveRunInvariant.append\_of\_keep} and
    \lname{Ratpoints.Sieve.SieveRunInvariant.quit\_eq\_and}, without
    \verb|SieveRunInvariant.append_of_forall|, the invariant of its loop with its stop
    (\lname{Ratpoints.Sieve.SieveRunLoopInvariant},
    \lname{Ratpoints.Sieve.SieveRunInvariant.eq\_deliver\_of\_quit}), and
    \lname{Ratpoints.Sieve.run\_sieveRun} for \verb|toList_run_sieveRun|, its contract the delivery
    of the list of before with the checks with the flags, and using Definition~\ref{def:deliver} and
    Lemma~\ref{lem:deliver} (G.6); the reduction of the run to the reference
    (Theorem~\ref{thm:run_reference}) names \lname{Ratpoints.Sieve.run\_sieveRun\_of\_eq\_filter}
    for \verb|toList_run_sieveRun_eq_flatMap| and \lname{List.flatMap\_filter\_flatMap\_filter}, its
    only user, delivers, uses Definition~\ref{def:deliver}, and its proof no longer uses
    Lemma~\ref{lem:flatMap_sound} (G.6); the points at infinity
    (Definition~\ref{def:pointsAtInfinity}) deliver call by call with the flags, without the test $s
    \ne 0$, name \lname{Ratpoints.Sieve.pointsAtInfinity} alone, give
    \csource{find_points.c}{5100}{5105} to the search and keep \csource{find_points.c}{5068}{5099},
    their claim about the candidate $(1, 0)$ of the reference holds for odd $d$ or $c_d \ne 0$, and
    their uses become Definitions~\ref{def:delivery} and~\ref{def:reversal} (G.7); the search for
    the points (Definition~\ref{def:findPoints}) is the instance of
    Definition~\ref{def:findPointsWith}, which takes all its links, and uses that node (G.7); the
    deviations after it lose the list of points for the callback and the work left to stage~G, gain
    the square root by \verb|Nat.sqrt| for \verb|mpz_sqrt| at infinity, and the sentence on the case
    $c_d = 0$ says that the model delivers $(1, 0, 0)$ twice there, as the C code would, which no
    input reaches (G.7); the sentence after the deviations of the search for the points says that
    the list holds the points without $\mathrm{noCheck}$, those with $y < 0$ left out under
    $\mathrm{noY}$, and the survivors under $\mathrm{noCheck}$ (G.8, Theorem~\ref{thm:noCheck} with
    G.9); Theorem~\ref{thm:findPoints} with the options of the output, as in
    Section~\ref{sec:outputOptions}, its proof using Lemmas~\ref{lem:findPoints_searchList}
    and~\ref{lem:noCheckList} (G.8); the paragraph after it names \verb|run_sieveRun| for
    \verb|toList_run_sieveRun| (G.6), states the list of Theorem~\ref{thm:findPoints_order} with
    $\mathrm{noCheck}$ false and filtered by $\mathrm{noY}$ (G.8), and gives the twelve statements
    of comparator and its record (G.14); the files section names \verb|Ratpoints/Proof/Output.lean|
    for the main theorems with their guards (G.7).
  \item Chapter~\ref{chap:words}: the check of a numerator (Definition~\ref{def:checkNum}) is the
    reference with the flag of the reversal, whose form with the flags is
    Definition~\ref{def:checkNumWith}; the walk, the scan, one piece and one denominator
    (Definitions~\ref{def:walkWord}, \ref{def:scanPiece}, \ref{def:siftPiece} and~\ref{def:sift})
    deliver and return at a stop, as in Section~\ref{sec:sieveDelivery}, the walk naming
    \lname{Ratpoints.Sieve.WalkDeliveryInvariant} and
    \lname{Ratpoints.Sieve.WalkDeliveryWhileInvariant} and using Definition~\ref{def:deliver},
    Lemmas~\ref{lem:checkDelivery}, \ref{lem:deliver} and~\ref{lem:sieveStops}, the scan naming
    \lname{Ratpoints.Sieve.ScanPieceInvariant.eq\_deliver\_of\_quit} and using
    Definition~\ref{def:deliver} and Lemmas~\ref{lem:deliver} and~\ref{lem:sieveStops}, one piece
    naming \lname{Ratpoints.Sieve.IsMasked.eq\_deliver} and using Definition~\ref{def:deliver}, and
    the node of one denominator naming
    \lname{Ratpoints.Sieve.flatMap\_acceptedPoints\_eq\_flatMap\_checkNumWith} for
    \verb|flatMap_acceptedPoints_eq_pointsOfDen| and using Definition~\ref{def:deliver} in place of
    Definition~\ref{def:checkNum}; Lemma~\ref{lem:sift_eq_checkDen} names
    \lname{Ratpoints.Sieve.run\_sift\_eq\_deliver\_checkDen} for \verb|toList_run_sift_eq_checkDen|,
    with the exact check on, uses Definition~\ref{def:deliver}, and its proof
    Lemma~\ref{lem:checkNumWith}; the paragraph after it says that the run delivers; the deviations
    lose the list of points for the callback and the work left to stage~G, and gain the returns at a
    stop (all G.6); the files section names \verb|Ratpoints/Proof/Scan.lean| for the walk and the
    scan (G.3).
  \item Chapter~\ref{chap:setup}: the introduction and the paragraph after
    Theorem~\ref{thm:refusals} say that membership holds with $\mathrm{noCheck}$ false, the set of
    elements with both options of the output off, and no duplicates for every option (G.8); the
    options (Definition~\ref{def:Options}) gain $\mathrm{noCheck}$ and $\mathrm{noY}$ and the link
    \csource{ratpoints.h}{403}{405}, and the sentence on the stages that add options says that
    stage~G added those of the output, which the reversal changes too (G.8); the coefficients as
    \verb|long|s (Definition~\ref{def:useCLong}) name \lname{Ratpoints.fitsLong\_natCast}, with the
    sentence that a natural number is a \verb|long| exactly when it is below $2^{63}$, which leaves
    the proof of Theorem~\ref{thm:refusals}, and Lemma~\ref{lem:reversal_swap} names
    \lname{Ratpoints.Setup.one\_le\_of\_reversal\_eq\_ok}, its statement gaining that $d' \ge 1$
    when $d \ge 1$, since the check lowers only an even $d \ge 2$, to $d - 1$; so the refusal
    theorems (Theorem~\ref{thm:refusals}), whose proofs now go through the search with a callback
    (\verb|exists_findPoints_eq_error_iff|), drop both names and use Theorem~\ref{thm:protocol} (all
    G.7); the map of the reversal (Definition~\ref{def:swapPoint}) is the one that the mapped call
    applies (G.1); Lemma~\ref{lem:swapPoint_bijection} names
    \lname{Ratpoints.Sieve.swapPoint\_swapPoint} in place of \verb|swapPoint_swapPoint_of_isPoint|,
    the map being its own inverse on the triples $(x, y, z)$ with $z > 0$ or $(x, z) = (1, 0)$,
    among them the normalized points, and \lname{Ratpoints.Sieve.swapPoint\_inj}, its injectivity on
    the same triples, and also \lname{Ratpoints.IsNormalizedPair} and
    \lname{Ratpoints.IsNormalizedPair.swapPoint}: a pair $(x, z)$ of integers is \emph{normalized of
    height at most $H$} when $\gcd(x, z) = 1$, $z > 0$ or ($z = 0$ and $x = 1$), $|x| \le H$ and $z
    \le H$, the conditions of a point (Definition~\ref{def:IsPoint}) without its equation, which the
    map of the reversal keeps, and \verb|isPoint_swapPoint| gets the conditions other than the
    equation from \verb|IsNormalizedPair.swapPoint| (G.8); the search
    (Definition~\ref{def:searchPoints}) delivers in the order of the C code, returns what a search
    gives back, names \lname{Ratpoints.searchPoints} alone, \verb|run_searchPoints| giving way to
    Lemma~\ref{lem:searchList}, which its node does not use, so that no cycle forms, loses the uses
    that \verb|run_searchPoints| had (Definition~\ref{def:searchDens}, Lemmas~\ref{lem:sievingInfo},
    \ref{lem:soundRun}, \ref{lem:survives} and~\ref{lem:twoAdic_exit}, Theorems~\ref{thm:bruteForce}
    and~\ref{thm:run_reference}), uses Definitions~\ref{def:sieving} and~\ref{def:Outcome}, loses
    \csource{find_points.c}{4573}{4573}, \csource{find_points.c}{5106}{5106},
    \csource{find_points.c}{5116}{5116}, \csource{find_points.c}{5119}{5119},
    \csource{find_points.c}{5132}{5132}, \csource{find_points.c}{5137}{5139} and
    \csource{find_points.c}{5167}{5167} (G.5) and gains \csource{find_points.c}{5100}{5106} and
    \csource{ratpoints.h}{415}{415} and \csource{ratpoints.h}{418}{418}, the bits
    \verb|RATPOINTS_CHECK_DENOM| and \verb|RATPOINTS_USE_JACOBI| that it sets (G.7); the paragraph
    after it describes $\mathrm{findPointsWith}$ (G.7); Theorem~\ref{thm:findPoints_perm} takes both
    options of the output off (G.8); the paragraph after Theorem~\ref{thm:refusals} on comparator
    (G.14); Section~\ref{sec:setupCli} lists only \verb|-l -u -dl -du| as not supported (G.11); the
    deviations lose the bits of the output left to stage~G and the list of points of the program
    (G.8 and G.11), and the deviation of the coefficients no longer says that the flags are not
    written back: the search writes back three bits (Definition~\ref{def:Outcome}), and the reversed
    coefficients and the degree stay unwritten (G.7); the sentence on the macros of the flags and of
    the return codes of \verb|ratpoints.h| says that stage~G records them, but for
    \verb|RATPOINTS_NO_REVERSE_AUTO|, which is stage~H's (G.14); the files section names
    \verb|Ratpoints/Proof/Scan.lean| for the walk and the scan (G.3) and
    \verb|Ratpoints/Proof/Output.lean| for the theorems (G.7).
  \item Chapter~\ref{chap:denominators}: the divisors of the leading coefficient
    (Definition~\ref{def:setupUs1}) find the primes in increasing order,
    \lname{Ratpoints.Setup.IsFactorization} with the ascending order, and their contracts
    \lname{Ratpoints.Setup.factorLcf\_spec} and \lname{Ratpoints.Setup.setupUs1\_spec} are one each:
    the cofactor $1$ exactly when the factorization is complete, then a factorization; and the
    result of the set-up for both outcomes, from those of the factorization and the divisors;
    \lname{Ratpoints.Setup.squarefreeDivisors\_spec} gives the fold of the doubling step or, once
    the loop has stopped, a prefix of it with at least $\mathrm{maxDivisors}$ entries, in place of
    the fold for a list of fewer than $\mathrm{maxDivisors}$ divisors, and the sentence of the node
    on it says so; the node names the invariant of the factorization
    (\lname{Ratpoints.Setup.CofactorInvariant},
    \lname{Ratpoints.Setup.FactorLcfInvariant.isRoughPart},
    \lname{Ratpoints.Setup.FactorLcfInvariant.iff\_and\_isFactorization}), the four facts on the
    table (\lname{Ratpoints.forIn\_toList\_primeTable},
    \lname{Ratpoints.two\_lt\_of\_mem\_toList\_primeTable},
    \lname{Ratpoints.eq\_of\_toList\_primeTable\_eq\_append\_singleton},
    \lname{Ratpoints.le\_of\_toList\_primeTable\_eq\_append\_cons\_cons}: it begins with $3$, ends
    with $1021$, holds only primes above $2$ and no prime between two neighbours) and the plain
    forms of the contracts (\lname{Ratpoints.Setup.run\_setupUs1\_eq},
    \lname{Ratpoints.Setup.snd\_run\_factorLcf\_eq\_one\_iff},
    \lname{Ratpoints.Setup.run\_squarefreeDivisors\_eq\_or}), and it uses
    Definition~\ref{def:binaryDivisors} and Lemma~\ref{lem:roughPart}; the list of the candidates
    for Mathlib names \verb|Ratpoints/Mathlib/NumberTheory/SmoothNumbers.lean| (all G.2); the
    paragraph after Lemma~\ref{lem:shape_dens} names the contract of the run on the delivery and its
    two equations (G.6); Section~\ref{sec:denSearch} names Lemma~\ref{lem:searchList} for
    \verb|run_searchPoints| and the C ranges of the search as they are now, with those of
    Definition~\ref{def:sieving} (G.5 and G.7); Theorem~\ref{thm:findPoints_order} takes
    $\mathrm{noCheck}$ false and filters by $\mathrm{noY}$,
    \lname{Ratpoints.eq\_bruteForce\_of\_findPoints\_eq\_ok\_of\_reverse\_eq\_false} and the
    permutation take both options of the output off (G.8), its list starts with what the reference
    finds over the candidate $(1, 0)$ (G.7), its proof uses Lemmas~\ref{lem:findPoints_searchList}
    and~\ref{lem:flatMap_sound} in place of Theorem~\ref{thm:refusals} (G.7), and a sentence refers
    to Theorem~\ref{thm:order} (G.10); the paragraph after it says that membership holds with
    $\mathrm{noCheck}$ false, the set of elements with both options of the output off and no
    duplicates for every option, and that the challenge's statement of membership has the options of
    the output (G.8), with the new statements of stage~G as they come (G.14); the files section
    names \verb|Ratpoints/Proof/Scan.lean| for the walk and the scan (G.3) and
    \verb|Ratpoints/Proof/Output.lean| for the theorems (G.7).
  \item Chapter~\ref{chap:tables}: the contract of the factorization of a modulus
    (Lemma~\ref{lem:modulusFactors}) uses Lemma~\ref{lem:roughPart}, for
    \lname{Array.pairwise\_map\_toList\_push}, which replaces the private lemma of
    \verb|Proof/Params.lean| in the steps of its invariant (G.2).
  \item Chapter~\ref{chap:stages}: the test of a candidate and the exact check
    (Definition~\ref{def:accepted}) deliver to the callback with the flags, and the node says that
    the points of the extraction for a numerator (\lname{Ratpoints.Sieve.acceptedPoints}) are those
    of the exact check with the flags (Definition~\ref{def:checkNumWith}) with the factor $b$ for
    odd $d$ and $1$ for even $d$ when the numerator passes the test of a candidate, and none
    otherwise, and uses Definition~\ref{def:checkNumWith}; the deviations gain the square root, and
    the test of the remainder, by \verb|Nat.sqrt| for \verb|mpz_sqrtrem| in the exact check;
    Lemma~\ref{lem:accepted} loses \verb|checkPoint_eq_checkNum|, whose place
    Lemma~\ref{lem:checkDelivery} takes, uses Definition~\ref{def:checkNumWith} in place of
    Definition~\ref{def:checkNum}, and \lname{Ratpoints.Sieve.acceptedPoints\_eq\_ite} gives the
    points of the extraction for a numerator $|a| \le H < 2^{63}$ over $0 < b < 2^{63}$ as those of
    the check with the flags when it survives the prime entries of the third stage, and none
    otherwise (G.6); a bit array through the second stage and the extraction
    (Definition~\ref{def:siftBitArray}) names \lname{Ratpoints.Sieve.siftBitArrayByWords} and
    \lname{Ratpoints.Sieve.siftBitArrayWhole} with their C ranges, and the deviation ``the arm of
    the second stage an argument, so that the guards run both in one build'' gains ``each arm a
    function of its own'' (G.4); Lemma~\ref{lem:siftBitArray} names
    \lname{Ratpoints.Sieve.siftBitArrayByWords\_spec} and
    \lname{Ratpoints.Sieve.siftBitArrayWhole\_spec} (G.4), and with G.6 the invariants of the stop
    (\lname{Ratpoints.Sieve.SiftBitArrayLoopInvariant},
    \lname{Ratpoints.Sieve.SiftBitArrayWalkLoopInvariant},
    \lname{Ratpoints.Sieve.SiftBitArrayInvariant.eq\_deliver\_of\_quit},
    \lname{Ratpoints.Sieve.SiftBitArrayWalkInvariant.eq\_deliver\_of\_quit},
    \lname{Ratpoints.Sieve.eq\_deliver\_of\_quit\_of\_forall\_eq}), using
    Definition~\ref{def:deliver} and, in its proof, Lemmas~\ref{lem:deliver}
    and~\ref{lem:sieveStops}; the deviations no longer leave the stop and the options of the output
    to stage~G (G.6 and G.8); the files section names \verb|Ratpoints/Proof/Scan.lean| (G.3) and
    \verb|Ratpoints/Proof/Output.lean| for the refusal theorems (G.7).
  \item Chapter~\ref{chap:packing}: the bits of one denominator (Definition~\ref{def:siftBits},
    Lemma~\ref{lem:siftBits}) deliver and return after a piece at which the callback stopped the
    search, the lemma naming \lname{Ratpoints.Sieve.SiftInvariant.eq\_deliver\_of\_quit}, using
    Definition~\ref{def:deliver} and, in its proof, Lemmas~\ref{lem:deliver}
    and~\ref{lem:sieveStops}; Section~\ref{sec:packedSieve} says that the contracts deliver (G.6);
    the files section names \verb|Ratpoints/Proof/Scan.lean| (G.3).
  \item Chapter~\ref{chap:rows}: the exclusions no longer leave the stop to stage~G (G.6).
  \item The stage list of the introduction refers to this chapter (Chapter~\ref{chap:output}), which
    is input after Chapter~\ref{chap:rows}. The generated table follows each item.
\end{itemize}

\section{Files, names and order of work}\label{sec:outputFiles}

The model, in the folder \verb|Ratpoints/Model|: the callback, the delivery, its reference, the
flags of the check and the appending callback in a new file without imports \verb|Delivery.lean|
(about $190$ lines with its guards); the exact check and the points at infinity in
\verb|Candidate.lean|, the sieve with the two arms in \verb|Sieve.lean|, the run in
\verb|Denominators.lean|, the options, what a search gives back, the sieving, the search, the search
with a callback and the search for the points in \verb|FindPoints.lean|, and the options of the
output, the callback, the formats and the report in \verb|Cli.lean| (about $940$ lines). The
specification: the order in a new file \verb|Ratpoints/Spec/Order.lean| (about $160$ lines). The
proofs, in the folder \verb|Ratpoints/Proof|, the contracts of the sieve still in files without
Mathlib: the reference of the delivery in a new core file \verb|Delivery.lean| and the count as a
sum in a new file \verb|DeliveryCount.lean|, with Mathlib; the points of a numerator and at infinity
with the flags in \verb|Stages.lean|; the walk, the arms and the scan in a new core file
\verb|Scan.lean|, split off \verb|Sift.lean| (about $970$ and $820$ lines); the run in
\verb|Denominators.lean|; the factorization and the divisors in \verb|Divisors.lean|; the points
that the search delivers, the sound tests and the survivors under \verb|RATPOINTS_NO_CHECK| in
\verb|FindPoints.lean|; the order in a new file \verb|Order.lean|; the search with a callback, the
protocol, the order theorem and the statements under \verb|RATPOINTS_NO_CHECK| in a new file
\verb|Output.lean| (about $520$ lines), with the main theorems, the theorems of the output of
stage~E and the refusal theorems, which leave \verb|FindPoints.lean| with their guards; the theorem
on the program in a new file \verb|Cli.lean|; normalized pairs in \verb|Setup.lean|; and the general
facts on core types in \verb|Containers.lean|, under their names in the namespaces of core. The
candidates for Mathlib: the rough part in a new mirror
\verb|Ratpoints/Mathlib/NumberTheory/SmoothNumbers.lean|, \lname{List.filter\_eq\_ite\_of\_forall}
and \lname{List.Nodup.flatMap} in the mirrors of \verb|Mathlib/Data/List|. Every Lean file stays
below $1500$ lines (the largest changed one, \verb|Ratpoints/Proof/FindPoints.lean|, about $1330$).
The new modules join their part roots and \verb|Ratpoints.lean|. One commit per item, each passing
the project's checks, the build with other values of the constants and the differential tests, with
the nodes and the passages of the earlier chapters that it changes and the C links moved from their
old owners in the same commit, and with the dependency audit of the blueprint after each item that
changes nodes; a lemma of a node comes with the item that first uses it, the items naming the lemmas
of a node that they bring and saying ``using'' for those that earlier items brought; the output of
\verb|ratpoints-lean| does not change up to G.10 and changes at G.11 and G.12:
\begin{enumerate}
  \item[G.1] The callback, the delivery, the flags of the check, the mapped call, the appending
    callback and the reference of the delivery (Definitions~\ref{def:delivery}
    and~\ref{def:deliver}), in \verb|Ratpoints/Model/Delivery.lean| with \verb|#guard| examples of
    stops at the $k$-th call with weights of $2$; their lemmas (Lemma~\ref{lem:deliver}) in
    \verb|Ratpoints/Proof/Delivery.lean| and \verb|Ratpoints/Proof/DeliveryCount.lean|;
    \verb|Answer| owns \csource{ratpoints.h}{436}{437} and \csource{ratpoints.h}{441}{442},
    \verb|CheckFlags| \csource{ratpoints.h}{414}{414}, each range later shared, as the record
    allows; the sentence of Definition~\ref{def:swapPoint} on the mapped call. Nothing uses them
    yet.
  \item[G.2] The factorization of the leading coefficient in increasing order and the outcomes of
    the divisor set-up (Lemmas~\ref{lem:roughPart} and~\ref{lem:binaryDivisors}), with the
    specification of the order, \verb|Ratpoints/Spec/Order.lean|, whose definitions the contracts
    read (Definitions~\ref{def:searchedPoly}, \ref{def:binaryDivisors} and~\ref{def:orderedPoints},
    with its examples), the mirror \verb|Ratpoints/Mathlib/NumberTheory/SmoothNumbers.lean|,
    \lname{Array.pairwise\_map\_toList\_push} in \verb|Proof/Containers.lean|, which replaces the
    private \verb|pairwise_map_fst_toList_push| of \verb|Proof/Params.lean| (two uses) and serves
    the new contract of the factorization in \verb|Proof/Divisors.lean|, and the four facts on the
    table in \verb|Params.lean| (\lname{Ratpoints.forIn\_toList\_primeTable},
    \lname{Ratpoints.two\_lt\_of\_mem\_toList\_primeTable},
    \lname{Ratpoints.eq\_of\_toList\_primeTable\_eq\_append\_singleton},
    \lname{Ratpoints.le\_of\_toList\_primeTable\_eq\_append\_cons\_cons}); in
    \verb|Ratpoints/Proof/Divisors.lean| one invariant and one contract of the factorization, one
    contract of the set-up for both outcomes, the contract of the doubling loop that stops, and
    \lname{Ratpoints.Setup.IsFactorization} with the ascending order;
    \verb|DivisorsInvariant.eq_foldl|, \verb|FactorLcfInvariant.isFactorization|,
    \verb|FactorLcfInvariant.suff_nil|, \verb|divisorsInvariant_of_le| and
    \verb|extendDivisorsInvariant_of_le| go; the three lemmas of the prototype on
    \lname{Ratpoints.FactorsByTrialDivision} in the namespace \verb|Nat|
    (\verb|Nat.IsRoughPart.factorsByTrialDivision_of_eq_one|,
    \verb|Nat.IsRoughPart.factorsByTrialDivision_of_lt|,
    \verb|Nat.IsRoughPart.eq_one_iff_factorsByTrialDivision|) move into \verb|Ratpoints| under names
    from their statements, the proposals
    \lname{Ratpoints.factorsByTrialDivision\_of\_isRoughPart\_one},
    \lname{Ratpoints.factorsByTrialDivision\_of\_isRoughPart\_of\_lt} and
    \lname{Ratpoints.eq\_one\_iff\_factorsByTrialDivision\_of\_isRoughPart}, which the prototype
    does not have yet; the node of Definition~\ref{def:setupUs1} (its names, its uses and its
    sentence on the doubling loop), the uses of Lemma~\ref{lem:modulusFactors}, the list of the
    candidates for Mathlib of Chapter~\ref{chap:denominators} and the files section of
    Chapter~\ref{chap:spec}. No model change; the challenge takes the import of the new file of the
    specification with G.10, when it restates it.
  \item[G.3] The walk, the extraction of one bit array and the scan of a piece moved from
    \verb|Ratpoints/Proof/Sift.lean| to a new core file \verb|Ratpoints/Proof/Scan.lean|, before the
    contracts on the delivery would take \verb|Sift.lean| past $1500$ lines; the type dump
    unchanged; the files sections of Chapters~\ref{chap:words}, \ref{chap:setup},
    \ref{chap:denominators}, \ref{chap:stages} and~\ref{chap:packing}.
  \item[G.4] The two arms of a bit array through the second stage and the extraction as functions of
    their own, since the stop tests of G.6 would take one function past $20$ lines, with their
    contracts in the form of today in \verb|Scan.lean|; each keeps the first numerator and the step
    where the C code computes them, as \verb|siftBitArray| does; the C ranges of the arms move from
    \verb|siftBitArray|, which keeps \csource{sift.c}{1094}{1111}, and its deviation gains ``each
    arm a function of its own''; the nodes of Definition~\ref{def:siftBitArray} and
    Lemma~\ref{lem:siftBitArray}, and the deviation of Chapter~\ref{chap:stages} on the arm of the
    second stage.
  \item[G.5] The sieving split off the search (Definition~\ref{def:sieving}), still on the list of
    the points; \csource{find_points.c}{4573}{4573}, \csource{find_points.c}{5116}{5116},
    \csource{find_points.c}{5119}{5119}, \csource{find_points.c}{5132}{5132},
    \csource{find_points.c}{5137}{5139} and \csource{find_points.c}{5167}{5167} move to it; the node
    of Definition~\ref{def:searchPoints} and Section~\ref{sec:denSearch}.
  \item[G.6] The sieve on the delivery, in one item: the exact check with the callback, the flags
    and the delivery, with its arms of \verb|RATPOINTS_NO_CHECK| and \verb|RATPOINTS_NO_Y|
    (Definition~\ref{def:checkNumWith} without \verb|pointsAtInfinityList|, the part of the check of
    Lemma~\ref{lem:checkNumWith}, \lname{Ratpoints.Sieve.checkNumWith\_of\_noCheck\_eq\_false},
    \lname{Ratpoints.Sieve.checkNumWith\_of\_noCheck\_eq\_true},
    \lname{Ratpoints.Sieve.checkNumWith\_eq\_map},
    \lname{Ratpoints.Sieve.mem\_checkNumWith\_of\_noCheck},
    \lname{Ratpoints.Sieve.map\_checkNumWith\_sublist},
    \lname{Ratpoints.Sieve.flatMap\_checkNumWith\_of\_noCheck\_eq\_false},
    \lname{Ratpoints.Sieve.fst\_snd\_swapPoint} and
    \lname{Ratpoints.Sieve.swapPoint\_fst\_zero\_snd}, and Lemma~\ref{lem:checkDelivery} without
    \verb|pointsAtInfinity_eq_deliver|), the walk, both arms of the extraction and their choice, the
    scan, one piece, the bits of one denominator, one denominator and the run, each returning at a
    stop, with their contracts on the delivery, the invariants of the stops in the nodes of the
    contracts and the prefixes of Lemma~\ref{lem:sieveStops}; \verb|checkPoint_eq_checkNum| and the
    helpers of the arrays of the walk, the arms, the scan, the masks, the bits and the run go,
    \verb|toList_run_sift_eq_checkDen|, \verb|flatMap_acceptedPoints_eq_pointsOfDen|,
    \verb|toList_run_sieveRun| and \verb|toList_run_sieveRun_eq_flatMap| give way to their forms on
    the delivery; the lines of the stops are already owned by these components, and their deviations
    lose the work left to stage~G. The sieving, still on the list of the points, calls the run with
    the appending callback on an empty array and the flag of the reversal alone, the one adapter of
    the item, which G.7 removes, so that the search and its proofs keep their statements; the flags
    of the check are those of the reversal until G.8. The adapter costs a corollary of
    \lname{Ratpoints.Sieve.run\_sieveRun\_of\_eq\_filter} for the appending callback on an empty
    array with the flag of the reversal alone, which gives the list of before, the reference over
    the denominators mapped by $\mathrm{swapPoint}$, through
    \lname{Ratpoints.Delivery.deliver\_appendPoint}, \lname{Ratpoints.Sieve.checkNumWith\_eq\_map},
    the prototype's \verb|Ratpoints.Sieve.checkNum_eq_checkNumWith| and
    \lname{Ratpoints.Sieve.pointsOfDen\_eq\_flatMap\_checkNum}, and the private lemmas of
    \verb|Proof/FindPoints.lean| on the run, one for each shape, restated through it; G.7 discards
    them. Comparator, since the proofs of the certified statements change. The nodes of
    Definitions~\ref{def:accepted}, \ref{def:walkWord}, \ref{def:siftBitArray}, \ref{def:scanPiece},
    \ref{def:siftPiece}, \ref{def:siftBits}, \ref{def:sift} and~\ref{def:sieveRun},
    Lemmas~\ref{lem:accepted}, \ref{lem:siftBitArray}, \ref{lem:siftBits}
    and~\ref{lem:sift_eq_checkDen} and Theorem~\ref{thm:run_reference}, the paragraph after
    Theorem~\ref{thm:findPoints} on \verb|run_sieveRun|, the paragraph after
    Lemma~\ref{lem:shape_dens}, Section~\ref{sec:packedSieve}, and the deviations and exclusions of
    Chapters~\ref{chap:words}, \ref{chap:stages} and~\ref{chap:rows} that leave the stop to stage~G.
  \item[G.7] The search on the delivery: the points at infinity call by call, without the test $s
    \ne 0$ (\verb|pointsAtInfinityList| of Definition~\ref{def:checkNumWith}, the part at infinity
    of Lemma~\ref{lem:checkNumWith},
    \lname{Ratpoints.Sieve.pointsAtInfinityList\_of\_noCheck\_eq\_false},
    \lname{Ratpoints.Sieve.pointsAtInfinityList\_of\_noCheck\_eq\_true},
    \lname{Ratpoints.Sieve.pointsAtInfinityList\_eq\_map},
    \lname{Ratpoints.Sieve.eq\_of\_mem\_pointsAtInfinityList\_of\_noCheck},
    \lname{Ratpoints.Sieve.map\_pointsAtInfinityList} and
    \lname{Ratpoints.Sieve.filter\_map\_swapPoint}, and
    \lname{Ratpoints.Sieve.pointsAtInfinity\_eq\_deliver} of Lemma~\ref{lem:checkDelivery}); the
    sieving on the delivery, the adapter, its corollary, the lemmas restated through it and the
    prototype's \verb|checkNum_eq_checkNumWith|, which nothing else uses, gone; the search in the
    order of the C code, with the return after the points at infinity; what a search gives back
    (Definition~\ref{def:Outcome}); the search with a callback (Definition~\ref{def:findPointsWith})
    and the search for the points as its instance with the appending callback, the flags of the
    check still those of the reversal; the points that the search delivers and the sound tests
    (Definition~\ref{def:searchList} and Lemma~\ref{lem:searchList}), \verb|run_searchPoints| giving
    way to its form with a sound test, \verb|pointsAtInfinity_eq_pointsOfDen| and
    \verb|pointsAtInfinity_eq_map| to the lemmas of the list at infinity; the new file
    \verb|Ratpoints/Proof/Output.lean| with the protocol (Theorem~\ref{thm:protocol}), the list of
    the search (Lemma~\ref{lem:findPoints_searchList}, its clause over the options of the output
    with G.8), and, moved there from \verb|Proof/FindPoints.lean| with their statements unchanged
    and their guards, the main theorems
    (\lname{Ratpoints.mem\_iff\_isPoint\_of\_findPoints\_eq\_ok},
    \lname{Ratpoints.nodup\_of\_findPoints\_eq\_ok},
    \lname{Ratpoints.toFinset\_eq\_points\_of\_findPoints\_eq\_ok}), the theorems of the output of
    stage~E (\lname{Ratpoints.exists\_eq\_of\_findPoints\_eq\_ok},
    \lname{Ratpoints.perm\_bruteForce\_of\_findPoints\_eq\_ok},
    \lname{Ratpoints.eq\_bruteForce\_of\_findPoints\_eq\_ok\_of\_reverse\_eq\_false}) and the
    refusal theorems (\lname{Ratpoints.exists\_findPoints\_eq\_error\_iff\_of\_sturm\_nonneg} and
    \lname{Ratpoints.exists\_findPoints\_eq\_error\_iff\_of\_sturm\_neg}), the file and module
    fields of \verb|formalization.yaml| following for the main theorems, the refusal theorems and
    the order; \lname{Ratpoints.fitsLong\_natCast} to Definition~\ref{def:useCLong} and
    \lname{Ratpoints.Setup.one\_le\_of\_reversal\_eq\_ok} to Lemma~\ref{lem:reversal_swap}, so that
    the refusal theorems use the protocol without a cycle; \csource{find_points.c}{5100}{5105} moves
    from the points at infinity to the search, the links of the search for the points to the search
    with a callback, which also owns \csource{ratpoints.h}{427}{430}, the prototypes and the mask of
    the bits of the input (\csource{ratpoints.h}{409}{411}); what a search gives back owns
    \csource{ratpoints.h}{377}{377}, \csource{ratpoints.h}{414}{414} and
    \csource{ratpoints.h}{416}{417}, the search \csource{ratpoints.h}{415}{415} and
    \csource{ratpoints.h}{418}{418}; \csource{ratpoints.h}{439}{439} (interface) and
    \csource{ratpoints.h}{444}{444} (memory) excluded; the challenge states the protocol, taking
    \verb|findPointsWith|, \verb|Answer|, \verb|Delivery| and \verb|Delivery.deliver| from the
    model, \verb|Solution| proves it, \verb|comparator/challenges.json| lists it, and comparator;
    the nodes of Definitions~\ref{def:pointsAtInfinity}, \ref{def:findPoints},
    \ref{def:searchPoints} and~\ref{def:useCLong}, Lemma~\ref{lem:reversal_swap},
    Theorems~\ref{thm:findPoints}, \ref{thm:findPoints_perm}, \ref{thm:refusals}
    and~\ref{thm:findPoints_order}, and Section~\ref{sec:denSearch}; the deviations of
    Chapter~\ref{chap:sieve} after Definition~\ref{def:findPoints} and that of the coefficients in
    Chapter~\ref{chap:setup}; the files sections of Chapters~\ref{chap:sieve}, \ref{chap:setup},
    \ref{chap:denominators} and~\ref{chap:stages}.
  \item[G.8] The options of the output (Section~\ref{sec:outputOptions}): the fields
    $\mathrm{noCheck}$ and $\mathrm{noY}$ of the options with \csource{ratpoints.h}{403}{405},
    copied into the flags of the check, and \csource{ratpoints.h}{406}{406} for the parameters;
    membership with $\mathrm{noCheck}$ false and the conclusion of $\mathrm{noY}$, no duplicates for
    every option, with normalized pairs and the survivors without duplicates
    (Lemma~\ref{lem:noCheckList} but for \verb|map_searchList_sublist|:
    \lname{Ratpoints.exists\_eq\_of\_mem\_searchList\_of\_noCheck},
    \lname{Ratpoints.nodup\_searchList\_of\_noCheck}, \lname{Ratpoints.searchList\_eq\_map},
    \lname{Ratpoints.nodup\_searchDens} and \lname{List.Nodup.flatMap}, using
    \verb|mem_checkNumWith_of_noCheck| and \verb|checkNumWith_eq_map| of G.6,
    \verb|eq_of_mem_pointsAtInfinityList_of_noCheck| and \verb|pointsAtInfinityList_eq_map| of G.7,
    and the range of the denominators of the search, which stays private in
    \verb|Proof/FindPoints.lean|), \lname{Ratpoints.Sieve.swapPoint\_swapPoint} for
    \verb|swapPoint_swapPoint_of_isPoint|, \lname{Ratpoints.Sieve.swapPoint\_inj},
    \lname{Ratpoints.IsNormalizedPair} and \lname{Ratpoints.IsNormalizedPair.swapPoint}, and the
    permutation, the set of elements and the reference itself with both options of the output off;
    the challenge's statement of membership with them, since \verb|Solution| would not build
    otherwise, and comparator; the nodes of Definitions~\ref{def:Options} and~\ref{def:Params},
    Lemma~\ref{lem:swapPoint_bijection} and Theorems~\ref{thm:findPoints}, \ref{thm:findPoints_perm}
    and~\ref{thm:findPoints_order}, and the passages that say that membership holds for every option
    (Section~\ref{sec:comparator-ab}, the sentence after the deviations of the search for the points
    and the paragraph after Theorem~\ref{thm:findPoints} in Chapter~\ref{chap:sieve}, the
    introduction and the paragraph after Theorem~\ref{thm:refusals} in Chapter~\ref{chap:setup}, and
    the paragraph after Theorem~\ref{thm:findPoints_order}), and the deviations of
    Chapters~\ref{chap:setup} and~\ref{chap:stages} that leave the options of the output to stage~G.
  \item[G.9] The statements under \verb|RATPOINTS_NO_CHECK| in the library
    (Theorem~\ref{thm:noCheck}), with the sublist of the list of the search
    (\lname{Ratpoints.map\_searchList\_sublist} of Lemma~\ref{lem:noCheckList}), using
    \verb|map_checkNumWith_sublist| of G.6 and \verb|map_pointsAtInfinityList| of G.7; the challenge
    states the three of them, \verb|Solution| proves them, \verb|comparator/challenges.json| lists
    them, and comparator; the reference to Theorem~\ref{thm:noCheck} in Chapter~\ref{chap:sieve}.
  \item[G.10] The order theorem (Theorem~\ref{thm:order}) with \verb|Ratpoints/Proof/Order.lean| and
    the lemmas of Lemma~\ref{lem:searchDenominators} in \verb|BruteForce.lean|, \verb|Shapes.lean|,
    \verb|Setup.lean|, \verb|FindPoints.lean| and the mirror of \verb|Mathlib/Data/List/Basic.lean|;
    the challenge restates \verb|Ratpoints/Spec/Order.lean| verbatim with its instance, imports
    \verb|Mathlib.Data.Nat.Factors|, its convention becoming the public imports of the specification
    files it restates, and states the order theorem, \verb|Solution| proves it,
    \verb|comparator/challenges.json| lists it, and comparator; the node of
    Theorem~\ref{thm:findPoints_order} and the paragraph after Definition~\ref{def:bruteForce}.
  \item[G.11] The program: the settings of the output, the options of the output read as
    \verb|main.c| reads them, the format, the strings and the callback of the program
    (Definitions~\ref{def:cliOutput}, \ref{def:formatPoint} and~\ref{def:cliProcess}), and the
    program on $\mathrm{findPointsWith}$, still without the report; the default format takes
    \csource{main.c}{681}{685} from the line of a point, which goes, the settings of the output own
    \csource{ratpoints.h}{407}{407}, the command line gains \csource{main.c}{533}{577},
    \csource{main.c}{602}{637} and \csource{main.c}{688}{688} and the head and the end of
    \verb|read_input| (\csource{main.c}{315}{316}, \csource{main.c}{691}{692}) for
    \verb|Ratpoints.Cli.input|, the reading of the coefficients takes their array
    (\csource{main.c}{54}{54}) for \verb|Ratpoints.Cli.coefficients|, the program the head and the
    end of \verb|main| (\csource{main.c}{195}{196}, \csource{main.c}{285}{285}), and the lines of
    \verb|main.c| of Section~\ref{sec:outputDeviations} but message 1 are excluded with their kinds;
    the harness compares the options of the output, the formats and the strings, \verb|-x| exactly
    against the pinned C program, its description of the pin saying that it puts the moduli of the
    model into the first two stages of both programs, and by the weaker check, the stop cases and
    the standard error, with its unit tests, a run without \verb|-q| still by its lines of points
    and pairs; the output changes: the options of the output take effect where they gave the status
    12; the nodes of Definitions~\ref{def:cliInput} and~\ref{def:cliRun}, Sections~\ref{sec:cli}
    and~\ref{sec:setupCli}, and the section on the differential tests of Chapter~\ref{chap:spec} but
    for the reports.
  \item[G.12] The report without \verb|-q| (Definition~\ref{def:cliReport}); the text of message 1
    (\csource{main.c}{977}{978}) excluded; the harness compares reports with the banner exchanged
    and the messages 3, 4 and 8 dropped on the side of the C program, the filtered runs and
    \verb|-v| by its lines; the output changes for the runs without \verb|-q|; the node of
    Definition~\ref{def:cliRun}, the list of Section~\ref{sec:cli} and the section on the
    differential tests of Chapter~\ref{chap:spec} for the reports.
  \item[G.13] The theorem on the program (Theorem~\ref{thm:cliRun}) in
    \verb|Ratpoints/Proof/Cli.lean|, with \lname{String.intercalate\_cons\_eq\_append\_join} in
    \verb|Containers.lean|.
  \item[G.14] Comparator with \verb|--paranoid| on the twelve statements and its record in
    \verb|comparator/README.md| and Chapter~\ref{chap:sieve}, with the texts on the certified
    statements (\verb|comparator/README.md|, Section~\ref{sec:comparator-ab}, the paragraphs after
    Theorems~\ref{thm:findPoints}, \ref{thm:refusals} and~\ref{thm:findPoints_order}); then the
    passages of the earlier chapters that G makes untrue (Section~\ref{sec:outputEarlier}) beyond
    those of the items, among them the sentence of Chapter~\ref{chap:setup} on the macros of
    \verb|ratpoints.h|, the description of the repository and \verb|formalization.yaml|.
\end{enumerate}
G.6 changes the sieve in one commit, since every split of it leaves an adapter in a caller that the
next split removes; its one adapter sits at its boundary, in the sieving, so that the search changes
once, in G.7.

\section{The evidence}\label{sec:output-evidence}

\begin{itemize}
  \item A reading of the C code of the callback, the stops, the options of the output and the flags
    written back, and of \verb|process|, the options, the formats, the strings and the report of
    \verb|main.c|, at the pin, with a second reader who checked it, on a build of the C program
    whose \verb|test4.sh| output is that of \verb|testbase4|: \verb|-1| stops at the first point
    that \verb|process| counts, also after a reversal, where the caller's points at infinity come
    from the sieve; message 4 after a stop at infinity counts no moduli; under \verb|-x| and
    \verb|-y| one point at infinity; the toggles, the escapes, a trailing \verb|%| and backslash,
    \verb|-f -x| taking \verb|-x| as the format, \verb|-ffoo| and a missing word as error 6, the
    strings around no points; \verb|-1|, \verb|-1 -i|, \verb|-x -1 -i| and \verb|-y -1 -i| give the
    same output with four threads; and no defect of the C code. The reading leaves notes on the C
    code, none a defect: the comment at \verb|main.c:951| lists \verb|\0|, which prints as it
    stands; \verb|%%| is no escape; a trailing \verb|%| or backslash skips the flush; and a sum of
    weights may collide with the codes of the errors.
  \item Before the stage, a sketch of the model in Lean in two parts, the library and the program,
    merged and checked against the C program: the default suites with every point list in its order
    and the reports compared whole, message 7 among them, but for the messages 3, 4 and 8, the $60$
    runs of \verb|test4.sh| with options of the output and no option of the search region, every one
    compared; $6536$ comparisons under \verb|-x| on $139$ curves of \verb|test4.sh| and $200$ random
    ones, exact with the moduli pinned and by the weaker check otherwise, which showed that the C
    program narrows the region of its Sturm step on a curve where the model does not, the reason for
    the comparison with \verb|-s|; the fuzz and the random curves with options of the output on two
    seeds; a callback that stops at its $k$-th call, with weights of either sign, on $416$ curves,
    which gets exactly the first $k$ points and their sum of weights; and the stop:
    \verb|'1 0 1' 1000000000000 -q -1| takes $7$~ms against $1$~ms for the C program, also at
    $2^{63} - 1$, and on the $41$ curves of \verb|test4.sh| whose first point lies in the sieve,
    \verb|-1| at the height $30000$ takes $5$ to $18$~ms against $20$ to $1591$~ms for the whole
    search, printing the first line of the whole output and of the C program; the cost of the
    listing of the denominators shows at large heights, $3.6$~seconds and $544$~MB for a stop in the
    sieve at $10^8$.
  \item Four reviews of the sketch, of the library against the C code, of the program, of the
    harness and of the order, each finding checked by a second reader, with runs of their own
    against the C program, none failing: for the library $4140$ comparisons on $138$ curves at the
    heights $7$, $60$ and $400$ under ten option sets with the stops, \verb|-y| and \verb|-x -s|
    with the moduli pinned, and message 7 the same on $114$ runs; for the program $16000$ runs
    compared byte for byte but for the messages 3, 4 and 8, so with message 7, on random option sets
    and curves; for the harness $2096$ runs under \verb|-x| compared exactly against the C program
    with \verb|-p 30 -N 30 -s|; and for the order $1535$ whole lists, $695$ more with \verb|-s|
    and $487$ runs with \verb|-y|. $44$ findings, $31$ confirmed, $10$ in part and $3$ refuted. They
    made the points at infinity follow the C code without the test $s \ne 0$, split the arms of the
    extraction, named the block of the sieving, stated the consequences of the deferred search
    region under \verb|RATPOINTS_NO_CHECK|, gave the harness the comparison of \verb|-x| against
    \verb|-p 30 -N 30 -s|, the empty standard error, the runtime that aborts on a panic and the stop
    cases, and made the factorization and the divisor set-up one contract each, with the rough part
    on Mathlib's smooth numbers.
  \item The definitions of the order checked against the C program line by line (\verb|-1 0 9| at
    the height $10$, \verb|0 1 0 2| at $10$, \verb|-6 5 -3 30| at $12$, \verb|1 0 0 30| at $20$,
    \verb|-2 0 0 1| at $130$, \verb|1 0 0 3| at $30$), three of them among the examples of the file.
  \item A prototype of all the proofs on the sketch: the whole library, the program, the challenge
    and the solution build, every proof at most $20$ lines with no new exception of the style lint,
    standard axioms; comparator with \verb|--paranoid| accepts the twelve statements with all six
    kernels, all of them stated at once in the prototype, where the order of work brings each with
    its theorem; the build with other values of the constants passes its nine variants,
    $\mathrm{maxDivisors}$ $2$ and $100000$ among them; the differential tests compare $3784$ times
    with none failed and $13$ times in \verb|test4high|; of the $281$ new public names, the six
    general ones are in neither Mathlib nor core nor Batteries. It changed the model of the sketch,
    as the reviews left it, in its docstrings only; the names of the three lemmas on
    $\mathrm{FactorsByTrialDivision}$ in G.2 are proposals of this chapter that the prototype does
    not have yet.
\end{itemize}
